\section{Maximal Ancestral Graphs}
\newcommand\NewEFCI[1]{{\color{orange}NewEFCI: #1}}

While the CDMG is already an abstraction of an SCM, we may wish to abstract the representation even further.
This will be especially useful for formulating constraint-based causal discovery.
A particular more abstract graphical representation that has seen much attention are \emph{maximal ancestral graphs (MAGs)} \cite{SMR1999,RichardsonSpirtes02}.
In order to deal with possible cycles (in simple SCMs), selection bias and exogenous input nodes, we extend the definition to what we will call $\sigma$-iMAGs.

\subsection{Additional graph terminology}

In this chapter, we will consider certain mixed graphs.
These mixed graphs can have two node types, where the default is `output' nodes, and optionally there can be `input' nodes (analogously to CDMGs).
The mixed graphs can have multiple edge types;
In addition to the three edge types for CDMGs ($\tuh$, $\hut$, $\huh$), there can be undirected edges ($\tut$).
Formally, we no longer introduce separate sets to represent the edges of each type, but merge them into the single set $E$.\footnote{So formally, every edge in $E$ is a tuple $(v_1,v_2,t)$ that consists of a pair of nodes $(v_1,v_2)$ and an edge type $t$, which is one of the four possibilities $\tuh$, $\hut$, $\huh$, $\tut$. We identify $(v_1,v_2,\tuh)\equiv (v_2,v_1,\hut)$, $(v_1,v_2,\tut)\equiv (v_2,v_1,\tut)$, $(v_1,v_2,\huh)\equiv (v_2,v_1,\huh)$.}
%Hence, a mixed graph is formally a tuple consisting of a set of nodes $V$ and a set of edges $E$ where each pair of nodes $(v_1,v_2)$ can be connected by a subset of the four edge types, i.e., $v_1 \tuh v_2$ and/or $v_1 \hut v_2$ and/or $v_1 \huh v_2$ and/or $v_1 \tut v_2$.}
Each edge $i \sus j$ has two \emph{edge marks}, one at each node, with each edge mark either a \emph{tail} or an \emph{arrowhead}. 
For example, the directed edge $i \tuh j$ has a tail at $i$ and an arrowhead at $j$, while the undirected edge $i \tut j$ has two tail edge marks.
All four possible combinations of edge marks can occur on an edge in a mixed graph. 
We will make use of the `$\asterisk$' symbol to denote any of the two edge marks. 
So the notation $i \sus j$ can stand for all four possible edge types between $i$ and $j$, whereas $i \hus j$ is shorthand for two possible edge types, as is $i \tus j$.
Edges of the form $i \hus j$ and $j \suh i$ are called \emph{into $i$}. 
Edges of the form $i \tus j$ and $j \sut i$ are called \emph{out of $i$}.

We extend various definitions of (directed) walks, (directed) paths and colliders to cover these new edge types in mixed graphs.
%\begin{Def}\label{def:mixed_graph_terminology}
%  Let $H=\lt J,V,E \rt$ be a mixed graph with input nodes $J$, output nodes $V$ and edges $E$ of the types $\{\tuh,\hut,\huh,\tut\}$.\footnote{Formally, we no longer introduce separate sets to represent the edges of each type, but merge them into the single set $E$.}
%  Let $v,w \in V \cup J$.
%  \begin{enumerate}
%    \item If there is an edge $v \sus w$ between $v$ and $w$ (of any type), we call $v$ and $w$ \emph{adjacent} in $H$.
%    \item Define $\Adj_H(v)$ to be the nodes adjacent to $v$ in $H$.
%    \item A triple of distinct nodes $\lt a,b,c \rt$ in $H$ form a \emph{triangle} if each pair of nodes in the triple is adjacent in $H$.
%    \item A triple of distinct nodes $\lt a,b,c \rt$ in $H$ is called \emph{unshielded} if $b$ is adjacent to both $a$ and $c$ in $H$, but $a$ is not adjacent to $c$ in $H$.
%    \item A \emph{walk between $v$ and $w$} in $H$ is a finite alternating sequence of nodes and edges 
%        \[v=v_0, a_0, v_1, \dots v_{n-1}, a_{n-1}, v_n=w\] 
%        in $H$ for some $n \ge 0$, i.e.\ such that for every $k=0,\dots,n-1$ 
%        we have that $v_{k},v_{k+1} \in V \cup J$ and $a_k = v_k \sus v_{k+1} \in E$, and with end nodes $v_0=v$ and $v_n=w$.
%    \item A walk is called a \emph{path} if no node occurs more than once on the walk.
%    \item A \emph{directed walk} (path) from $v$ to $w$ in $H$ is a walk (path) of the form:
%        \[v=v_0 \tuh v_1 \tuh  \cdots \tuh v_{n-1} \tuh v_n=w,\]
%        for some $n \ge 0$.
%    \item A \emph{directed cycle} is a directed walk $v \tuh \dots \tuh w$, concatenated with the directed edge $w \tuh v$.
%    \item An \emph{almost directed cycle} is a directed walk $v \tuh \dots \tuh w$, concatenated with the bidirected edge $w \huh v$.
%    \item A path $v_0 \sus \dots \sus v_n$ in $H$ is called \emph{uncovered} if every subsequent triple $\lt v_{k-1}, v_k, v_{k+1} \rt$ for $1 < k < n$ is unshielded.
%    \item A triple of consecutive nodes $v_{k-1} \sus v_k \sus v_{k+1}$ on a walk in $H$ is called \emph{collider} if it is of the form $v_{k-1} \suh v_k \hus v_{k+1}$.
%    \item A triple of consecutive nodes $v_{k-1} \sus v_k \sus v_{k+1}$ on a walk in $H$ is called a \emph{non-collider} if it is of the form $v_{k-1} \sus v_k \tus v_{k+1}$ or $v_{k-1} \sut v_k \sus v_{k+1}$.  Furthermore, we also refer to the end nodes $v_0$ and $v_n$ of a walk between $v_0$ and $v_n$ in $H$ as \emph{non-colliders}.
%    \item A path between $v$ and $w$ in $H$ is called a \emph{collider path} if every non-endpoint node on the path is a collider on the path.
%    \item If there is a directed walk from $v$ to $w$ in $H$ then we say that \emph{$v$ is ancestor of $w$ in $H$}, and we write $v \in \Anc_H(w)$. For $W \subseteq V \cup J$, we define $\Anc_H(W) := \bigcup_{w\in W} \Anc_H(w)$.
%    \item If there is a directed walk from $v$ to $w$ in $H$ then we say that \emph{$w$ is descendant of $v$ in $H$}, and we write $w \in \Desc_H(v)$. For $W \subseteq V \cup J$, we define $\Desc_H(W) := \bigcup_{w\in W} \Desc_H(w)$.
%    \item A walk between $v,w \in V \cup J$ (with $v \ne w$) in $H$ is called \emph{inducing} if every non-endpoint node is a collider in $\Anc_H(\{v,w\})$.
%    \item If $v \tut w$ in $H$ then $v,w$ are called \emph{neighbors} in $H$.
%    \item The set of neighbors of a node $v$ in $H$ is denoted $\Nb^H(v) := \{w : v \tut w \in H\}$ and called the \emph{neighborhood of $v$}; 
%      for a subset $A$ of nodes in $H$, we define $\Nb^H(A) := \bigcup_{v \in A} \Nb^H(v)$.
%      If for each pair $w_1,w_2 \in \Nb^H(v)$ there is an edge $w_1 \tut w_2$ in $H$ provided $w_1 \ne w_2$, we say that the neighborhood of $v$ is \emph{complete}.
%  \end{enumerate}
%\end{Def}
\begin{Def}\label{def:mixed_graph_terminology_adjacency}
  \Joris{These are literally the same for PAGs.}
  Let $H=\lt J,V,E \rt$ be a mixed graph with input nodes $J$, output nodes $V$ and edges $E$ of the types $\{\tuh,\hut,\huh,\tut\}$.
  %Let $H=\lt J,V,E \rt$ be a mixed graph with input nodes $J$, output nodes $V$ and edges $E$ of the types $\{\tuh,\hut,\huo,\huh,\ouo,\ouh,\tuo,\out,\tut\}$.
  Let $v,w \in V \cup J$.
  \begin{enumerate}
    \item If there is an edge $v \sus w$ between $v$ and $w$ (of any type), we call $v$ and $w$ \emph{adjacent} in $H$.
    \item Define $\Adj_H(v)$ to be the nodes adjacent to $v$ in $H$.
    \item A triple of distinct nodes $\lt a,b,c \rt$ in $H$ form a \emph{triangle} if each pair of nodes in the triple is adjacent in $H$.
    \item A triple of distinct nodes $\lt a,b,c \rt$ in $H$ is called \emph{unshielded} if $b$ is adjacent to both $a$ and $c$ in $H$, but $a$ is not adjacent to $c$ in $H$.
    \item A \emph{walk between $v$ and $w$} in $H$ is a finite alternating sequence of nodes and edges 
        \[v=v_0, a_0, v_1, \dots v_{n-1}, a_{n-1}, v_n=w\] 
        in $H$ for some $n \ge 0$, i.e.\ such that for every $k=0,\dots,n-1$ 
        we have that $v_{k},v_{k+1} \in V \cup J$ and $a_k = v_k \sus v_{k+1} \in E$, and with end nodes $v_0=v$ and $v_n=w$.
    \item A walk is called a \emph{path} if no node occurs more than once on the walk.
    \item A path $v_0 \sus \dots \sus v_n$ in $H$ is called \emph{uncovered} if every subsequent triple $\lt v_{k-1}, v_k, v_{k+1} \rt$ for $1 < k < n$ is unshielded.
  \end{enumerate}
\end{Def}
\begin{Def}\label{def:mixed_graph_terminology_orientations}
  \Joris{These translate into `definite' and `possible' notions in PAGs.}
  Let $H=\lt J,V,E \rt$ be a mixed graph with input nodes $J$, output nodes $V$ and edges $E$ of the types $\{\tuh,\hut,\huh,\tut\}$.
  Let $v,w \in V \cup J$.
  \begin{enumerate}
    \item A \emph{directed walk} (path) from $v$ to $w$ in $H$ is a walk (path) of the form:
        \[v=v_0 \tuh v_1 \tuh  \cdots \tuh v_{n-1} \tuh v_n=w,\]
        for some $n \ge 0$.
    \item A \emph{directed cycle} is a directed walk $v \tuh \dots \tuh w$, concatenated with the directed edge $w \tuh v$.
    \item An \emph{almost directed cycle} is a directed walk $v \tuh \dots \tuh w$, concatenated with the bidirected edge $w \huh v$.
    \item If there is a directed walk from $v$ to $w$ in $H$ then we say that \emph{$v$ is ancestor of $w$ in $H$}, and we write $v \in \Anc_H(w)$. For $W \subseteq V \cup J$, we define $\Anc_H(W) := \bigcup_{w\in W} \Anc_H(w)$.
    \item If there is a directed walk from $v$ to $w$ in $H$ then we say that \emph{$w$ is descendant of $v$ in $H$}, and we write $w \in \Desc_H(v)$. For $W \subseteq V \cup J$, we define $\Desc_H(W) := \bigcup_{w\in W} \Desc_H(w)$.
    \item A triple of consecutive nodes $v_{k-1} \sus v_k \sus v_{k+1}$ on a walk in $H$ is called \emph{collider} if it is of the form $v_{k-1} \suh v_k \hus v_{k+1}$.
    \item A triple of consecutive nodes $v_{k-1} \sus v_k \sus v_{k+1}$ on a walk in $H$ is called a \emph{non-collider} if it is of the form $v_{k-1} \sus v_k \tus v_{k+1}$ or $v_{k-1} \sut v_k \sus v_{k+1}$. 
      Furthermore, we also refer to the end nodes $v_0$ and $v_n$ of a walk between $v_0$ and $v_n$ in $H$ as \emph{non-colliders}.
    \item A path between $v$ and $w$ in $H$ is called a \emph{collider path} if every non-endpoint node on the path is a collider on the path.
    \item A walk between $v,w \in V \cup J$ (with $v \ne w$) in $H$ is called \emph{inducing} if every non-endpoint node is a collider in $\Anc_H(\{v,w\})$.
  \end{enumerate}
  \Joris{Some recently added ones:}
  \begin{enumerate}
    \item If $v \tut w$ in $H$ then $v,w$ are called \emph{neighbors} in $H$.
    \item The set of neighbors of a node $v$ in $H$ is denoted $\Nb^H(v) := \{w : v \tut w \in H\}$ and called the \emph{neighborhood of $v$}; 
      for a subset $A$ of nodes in $H$, we define $\Nb^H(A) := \bigcup_{v \in A} \Nb^H(v)$.
      If for each pair $w_1,w_2 \in \Nb^H(v)$ there is an edge $w_1 \tut w_2$ in $H$ provided $w_1 \ne w_2$, we say that the neighborhood of $v$ is \emph{complete}.
    \item An \emph{anterior walk} (path) from $v$ to $w$ in $H$ is a walk (path) of the form:
        \[v=v_0 \tus v_1 \tus \cdots \tus v_{n-1} \tus v_n=w,\]
        for some $n \ge 0$.
    \item If there exists an anterior walk from $v$ to $w$ in $H$ then we say that \emph{$v$ is anterior to $w$ in $H$}, and we write $v \in \Ant^H(w)$. For $W \subseteq V \cup J$, we define $\Ant^H(W) := \bigcup_{w\in W} \Ant^H(w)$.
  \end{enumerate}
\end{Def}

\subsection{Acyclic setting}

MAGs were introduced in \cite{RichardsonSpirtes02}; we have slightly modified the definition to allow for input nodes, and refer to these objects as `iMAGs' to avoid possible confusion with MAGs.
\begin{Def}\label{def:imag}
  A mixed graph $H = \lt J,V,E\rt$ with input nodes $J$, output nodes $V$ and edges $E$ of the types $\{\tuh,\hut,\huh,\tut\}$ is called a \emph{maximal ancestral graph with input nodes (iMAG)} if all of the following conditions hold:
  \begin{enumerate}
    \item Between any two distinct nodes there is at most one edge, and there are no edges between a node and itself;
    \item The pattern $a \suh b \tut c$ (an edge into an undirected edge) does not occur in $H$;
    \item If $H$ contains an anterior path $a \tus \dots \tus b$ (for $a \ne b$) then it does not contain an edge $a \hus b$; (`ancestral');
    \item There is no inducing path between any two distinct non-adjacent nodes (`maximal');
    \item Only undirected edges are allowed between pairs of input nodes (if $j_1,j_2 \in J$ are adjacent in $H$, then $j_1 \tut j_2$ in $H$); 
      furthermore, an edge between an input node $j \in J$ and an output node $v \in V$ must be of the form $j \tus v$.
  \end{enumerate}
  We call it a \emph{maximal ancestral graph (MAG)} if it has no input nodes (i.e., $J=\emptyset$).
\end{Def}
\Joris{Note that alternatively, we might reorganize this definition by first defining MAGs (without inputs), and then defining
an iMAG as a MAG with nodes $J \dcup V$ such that the last item in Definition~\ref{def:imag} holds (in addition to the other 
points, which don't distinguish between the two node types and hold when thinking about the iMAG just as a MAG). This would
show very clearly that all properties that hold for MAGs are automatically true for iMAGs, but iMAGs have
some additional constraints regarding the distinction of input/output nodes. In other words, iMAGs behave like MAGs when
ignoring the difference between input/output nodes, but have some additional properties that concern the differences of the two nodes types.}
We use iMAGs to represent sets of CADMGs in the following sense.
\begin{Def}\label{def:imag_represents_cadmg}
  Let $H = \lt J,V,E\rt$ be a mixed graph with input nodes $J$, output nodes $V$ and edges $E$ of the types $\{\tuh,\hut,\huh,\tut\}$.
  Let $G$ be a CADMG with input nodes $J$ and output nodes $V^+ = V \dcup S$.
  We say that \emph{$H$ represents $G$ given $S$} if all of the following hold:
  \begin{enumerate}
    \item Between any two distinct nodes in $H$ there is at most one edge, and there are no edges between a node and itself;
    \item Two distinct nodes $i,j \in V \cup J$ are adjacent in $H$ if and only if there is a inducing path between $i$ and $j$ given $S$ in $G$;
    \item If $i \hus j$ in $H$ then $i \notin \Anc_{G}(\{j\} \cup S)$;
    \item If $i \tus j$ in $H$ then $i \in \Anc_{G}(\{j\} \cup S)$.
  \end{enumerate}
\end{Def}
If input nodes are absent ($J=\emptyset$) this specializes to how MAGs represent ADMGs.
We conjecture that the abstract Definition~\ref{def:imag} provides an axiomatic characterization of the class of mixed graphs that represent a CADMG:
\begin{Con}\label{con:imags_represent_cadmgs}
  Let $H = \lt J,V,E\rt$ be a mixed graph with input nodes $J$, output nodes $V$ and edges $E$ of the types $\{\tuh,\hut,\huh,\tut\}$.
  The following equivalence holds:  
  $H$ is an iMAG $\iff$ there exists a CADMG $G$ with input nodes $J$ and output nodes $V^+ = V \dcup S$ such that $H$ is a mixed graph that represents $G$ given $S$.
\end{Con}
\begin{proof}
  \Joris{I don't remember having seen a proof for this statement for the case $J=\emptyset$, but it must be true because that's precisely how MAGs are used to represent ADMGs and what must have inspired the axiomatic definition of MAGs in the first place.

  The next bit considers what happens in case $J\ne \emptyset$:

  If $H$ represents $G$ given $S$ then if there is an edge between input nodes $j_1, j_2$ then there must be a $d$-inducing path in $G$ between $i$ and $j$ given $S$.
  Every non-endpoint node must be a collider that is ancestor of $S$ (because it cannot be ancestor of $j_1$ or $j_2$ by assumption of the class of CADMGs).
  It must have at least one collider because otherwise $j_1,j_2$ would be adjacent in the CADMG.
  Hence, on that path, $j_1$ must be ancestor of $S$, and $j_2$ must be ancestor of $S$. 
  Hence we end up with an undirected edge between $j_1$ and $j_2$ in $H$!
  Can we ever end up with an arrowhead at $j$ in $H$?
  Not on an edge connecting it with another input node as we have just seen.
  So consider an edge $j \sus v$ in $H$.
  There must be a $d$-inducing path between $j$ and $v$ given $S$ in $G$ of the form $j \tuh ... \sus v$.
  This path either continues as a directed path the whole way (in which case $j$ is ancestor of $v$ in $G$) or it must meet a collider in between.
  The collider must be ancestor of $v \cup S$.
  Hence $j$ must be ancestor of $v \cup S$.
  Hence we must have a tail $j \tus v$ in $H$.
  So this together gives one direction of the equivalence.
  The other still needs to be written down (by explicit construction of such a $G$).}
\end{proof}

\subsection{Abstract definition of $\sigma$-iMAGs}

We will now consider the cyclic setting so that we can represent CDMGs corresponding to simple SCMs.
This combines three extensions of the usual definition of MAG: we allow for input nodes, 
there can now be edges into undirected edges (under certain conditions), and we have weakened the 
condition of being ancestral to $\sigma$-ancestral.
\begin{Def}\label{def:sigmaimag}
  A mixed graph $H = \lt J,V,E\rt$ with input nodes $J$, output nodes $V$ and edges $E$ of the types $\{\tuh,\hut,\huh,\tut\}$ is called a \emph{$\sigma$-maximal ancestral graph with input nodes ($\sigma$-iMAG)} if all of the following conditions hold:
  \begin{enumerate}
    \item Between any two distinct nodes there is at most one edge, and there are no edges between a node and itself;
    \item If $H$ contains an anterior path $a \tus \dots \tus b$ (for $a \ne b$) then it does not contain an edge $a \hus b$; (`ancestral');
    \item $H$ is `$\sigma$-maximal':
      \begin{enumerate}
        \item There is no inducing path between any two distinct non-adjacent nodes (`maximal'), and
        \item If $a \suh b \tut c$ is in $H$ then $a$ and $c$ must be adjacent in $H$;
          if furthermore $b \tut d$ is in $H$ then $b$ and $d$ must also be adjacent in $H$ (`$\sigma$-complete').
          \Joris{This somewhat weird looking axiom might be due to our definition of `inducing path'; perhaps if 
          we were to also consider $a \suh b \tut c$ to be inducing paths we could drop this second point 
          because it would already be implied by the first one (`maximal'); however, the notion of `inducing path'
          needs to satisfy various other `natural' properties as well, so if one starts to extend the definition
          to also consider $a \suh b \tut c$ one should think which other paths should then also be considered
          inducing in order to preserve/obtain the nice properties. Essentially the goal should be to find the
          `right' definition of inducing path that leads to the simplest proofs overall... This might actually
          also involve redefining the notion of $\sigma$-inducing path to make the two notions nicely compatible...
          it's not clear currently whether this would bring any advantages}
      \end{enumerate}
    \item Only undirected edges are allowed between pairs of input nodes (if $j_1,j_2 \in J$ are adjacent in $H$, then $j_1 \tut j_2$ in $H$); 
      furthermore, an edge between an input node $j \in J$ and an output node $v \in V$ must be of the form $j \tus v$.
  \end{enumerate}
  We call it a \emph{$\sigma$-maximal ancestral graph ($\sigma$-MAG)} if it has no input nodes (i.e., $J=\emptyset$).
\end{Def}
\Joris{Note that alternatively, we might reorganize this definition by first defining $\sigma$-MAGs (without inputs), and then defining
a $\sigma$-iMAG as a $\sigma$-MAG with nodes $J \dcup V$ such that the last item in Definition~\ref{def:sigmaimag} holds. This would
show very clearly that all properties that hold for $\sigma$-MAGs are automatically true for $\sigma$-iMAGs, but $\sigma$-iMAGs have
some additional properties regarding the distinction of input/output nodes.}
We use $\sigma$-iMAGs to represent sets of CDMGs in the following sense.
\begin{Def}\label{def:sigmaimag_represents_cdmg}
  Let $H = \lt J,V,E\rt$ be a mixed graph with input nodes $J$, output nodes $V$ and edges $E$ of the types $\{\tuh,\hut,\huh,\tut\}$.
  Let $G$ be a CDMG with input nodes $J$ and output nodes $V^+ = V \dcup S$.
  We say that \emph{$H$ represents $G$ given $S$} if all of the following hold:
  \begin{enumerate}
    \item Between any two distinct nodes in $H$ there is at most one edge, and there are no edges between a node and itself;
    \item Two distinct nodes $i,j \in V \cup J$ are adjacent in $H$ if and only if there is a $\sigma$-inducing path between $i$ and $j$ given $S$ in $G$;
    \item If $i \hus j$ in $H$ then $i \notin \Anc_{G}(\{j\} \cup S)$;
    \item If $i \tus j$ in $H$ then $i \in \Anc_{G}(\{j\} \cup S)$.
  \end{enumerate}
\end{Def}
This is almost literally the same as Definition~\ref{def:imag_represents_cadmg}, except that `inducing path' is replaced by `$\sigma$-inducing path'.
If input nodes are absent ($J=\emptyset$) this specializes to how $\sigma$-MAGs represent DMGs.
We conjecture that the abstract Definition~\ref{def:sigmaimag} provides an axiomatic characterization of the class of mixed graphs that represent a CDMG:
\begin{Con}\label{con:sigmaimags_represent_cdmgs}
  Let $H = \lt J,V,E\rt$ be a mixed graph with input nodes $J$, output nodes $V$ and edges $E$ of the types $\{\tuh,\hut,\huh,\tut\}$.
  The following equivalence holds:  
  $H$ is a $\sigma$-iMAG $\iff$ there exists a CDMG $G$ with input nodes $J$ and output nodes $V^+ = V \dcup S$ such that $H$ is a mixed graph that represents $G$ given $S$.
\end{Con}
\begin{proof}
  \Joris{Binghua and I proved the special case $J=\emptyset$.

  The next bit considers what happens in case $J\ne \emptyset$:
  
  If $H$ represents $G$ given $S$ then if there is an edge between input nodes $j_1, j_2$ then there must be a $\sigma$-inducing path in $G$ between $i$ and $j$ given $S$.
  Every non-endpoint node must either be an unblockable non-collider, or a collider that is ancestor of $S$ (because it cannot be ancestor of $j_1$ or $j_2$ by assumption of the class of DMGs).
  It must have at least one collider because otherwise $j_1,j_2$ end up in the same scc.
  Hence, on that path, $j_1$ must be ancestor of $S$, and $j_2$ must be ancestor of $S$. 
  Hence we end up with an undirected edge between $j_1$ and $j_2$ in $H$!
  Can we ever end up with an arrowhead at $j$ in $H$?
  Not on an edge connecting it with another input node as we have just seen.
  So consider an edge $j \sus v$ in $H$.
  There must be a $\sigma$-inducing path between $j$ and $v$ given $S$ in $G$ of the form $j \tuh ... \sus v$.
  This path either continues as a directed path the whole way (in which case $j$ is ancestor of $v$ in $G$) or it must meet a collider in between.
  The collider must be ancestor of $v \cup S$.
  Hence $j$ must be ancestor of $v \cup S$.
  Hence we must have a tail $j \tus v$ in $H$.

  So this together gives one direction of the equivalence.
  The other still needs to be written down (by explicit construction of such a $G$).
  I would just hope that the construction designed for the case $J=\emptyset$ automatically works as well for $J\ne \emptyset$?}
\end{proof}


\Joris{
  The current setting appears to be `internally consistent', that is, because we have chosen (variation) independent input nodes in the CDMG setting,
  we never need edges between input nodes for CDMGs, and because we assume that the exogenous random variables are `randomized' (that is, probabilistically
  independent) w.r.t.\ the exogenous inputs, and the endogenous variables can never cause the exogenous inputs, hence there can be only outgoing directed
  edges from input nodes to output nodes in CDMGs, we end up working with this specific class of ($\sigma$-)iMAGs.

  We might consider relaxing the framework at the CDMG level, in two ways:
  \begin{enumerate}
    \item We could allow CDMGs to have undirected edges between inputs (to model variation dependencies between inputs stemming from a
      selection mechanism that is chronologically \emph{after} the interventions modeled by the input nodes but \emph{before} the possible
      interventions on endogenous variables. I would guess that this does \emph{not} require changing the definition of ($\sigma$-)iMAGs.
      Actually, now that I write this down, it seems a very sensible/natural way to define the conditioning operation on the CDMG level!
    \item We could further relax by allowing CDMGs to have any type of edges (undirected, directed, bidirected) between inputs, and even 
      multiple ones. The idea would be that we may have background knowledge about the causal relationships between the input variables.
      One may ask whether it would be \emph{useful} to do this because one may believe that this is `overly specific' modeling: the 
      precise edge types between input nodes may have no consequences down the road, and WLOG one could replace them by only one edge
      type (e.g., undirected). So currently this seems not of high priority; but we can still ask: if we would allow this, would this just lead
      us to drop the point 5 in Definition~\ref{def:imag} and point 4 in Definition~\ref{def:sigmaimag}, that is, they just don't need to 
      distinguish between input/output nodes anymore?
  \end{enumerate}
}

\subsubsection{Obsolete?}

\begin{Def}\label{def:sigmamag_alt2}
  A mixed graph $H = \lt J,V,E\rt$ with input nodes $J$, output nodes $V$ and edges $E$ of the types $\{\tuh,\hut,\huh,\tut\}$ is called a \emph{maximal $\sigma$-ancestral graph with input nodes ($\sigma$-iMAG)} if all of the following conditions hold:
  \begin{enumerate}
    \item Between any two distinct nodes there is at most one edge, and there are no edges between a node and itself;
    \item If $a \suh b$ then $\Nb^H(b)$ must be complete; also, $H$ contains no unshielded triples of the form $a \suh b \tut c$ (`$\sigma$-complete');
%    \item $H$ contains no directed cycles and no almost directed cycles; %, and no unshielded triples $a \suh b \tut c$;
%      if $a \tuh b \tut c$ in $H$, then $a \tuh c$ in $H$, and if $a \huh b \tut c$ in $H$ then $a \huh c$ in $H$ (`$\sigma$-ancestral');
    \item If $H$ contains an anterior path $a \tus \dots \tus b$ (for $a \ne b$) then it does not contain an edge $a \hus b$; (`$\sigma$-ancestral');
    \item There is no inducing path between any two distinct non-adjacent nodes (`maximal');
    \item No input node is adjacent to any other input node. \Joris{perhaps we should just drop this one and replace it with: only undirected edges allowed between input nodes; or even that is not general enough?}
  \end{enumerate}
  We call it a \emph{maximal $\sigma$-ancestral graph ($\sigma$-MAG)} if it has no input nodes (i.e., $J=\emptyset$).
\end{Def}

An alternave (less generalizable) definition would be:
\begin{Def}\label{def:sigmamag_alt}
  A mixed graph $H = \lt J,V,E\rt$ with input nodes $J$, output nodes $V$ and edges $E$ of the types $\{\tuh,\hut,\huh,\tut\}$ is called a \emph{maximal $\sigma$-ancestral graph with input nodes ($\sigma$-iMAG)} if all of the following conditions hold:
  \begin{enumerate}
    \item Between any two distinct nodes there is at most one edge, and there are no edges between a node and itself;
%    \item If $a \suh b$ then $\Nb^H(b)$ must be complete; also, $H$ contains no unshielded triples of the form $a \suh b \tut c$ (`$\sigma$-complete');
    \item $H$ contains no directed cycles and no almost directed cycles; %, and no unshielded triples $a \suh b \tut c$;
      if $a \tuh b \tut c$ in $H$, then $a \tuh c$ in $H$, and if $a \huh b \tut c$ in $H$ then $a \huh c$ in $H$ (`$\sigma$-ancestral');
%    \item If $H$ contains an anterior path $a \tus \dots \tus b$ (for $a \ne b$) then it does not contain an edge $a \hus b$; (`$\sigma$-ancestral');
    \item There is no inducing path between any two distinct non-adjacent nodes (`maximal');
    \item No input node is adjacent to any other input node. \Joris{perhaps we should just drop this one and replace it with: only undirected edges allowed between input nodes; or even that is not general enough?}
  \end{enumerate}
  We call it a \emph{maximal $\sigma$-ancestral graph ($\sigma$-MAG)} if it has no input nodes (i.e., $J=\emptyset$).
\end{Def}

The following states that there are no `(almost) anterior cycles' in a $\sigma$-iMAG, analogously to how there are no (almost) directed cycles.
This shows the equivalence of both definitions.
\begin{Lem}\label{lem:sigmamag_no_(almost)_anterior_path_cycles}
  Let $H$ be a $\sigma$-iMAG with input nodes $V$ and output nodes $J$.
  Let $a, b \in V \cup J$ be two distinct nodes.
  If there is an anterior path from $a$ to $b$ in $H$, there cannot be an edge $a \hus b$ in $H$.
\end{Lem}
\begin{proof}
  Let $\pi = (x_0 \tus x_1 \tus \dots \tus x_n)$ be a shortest anterior path from $a=x_0$ to $b=x_n$ in $H$.
  We must have $n > 0$ because $a \ne b$, and $n \ne 1$ because otherwise we would have both $a \tus b$ and $a \hus b$ in $H$.
  Hence $n \ge 2$.

  Suppose an edge $a \hus b$ is in $H$.

  First assume $\pi$ is undirected.
  Then we have $b \suh a \tut x_1$ in $H$, and hence also $b \suh x_1$ in $H$.
  Furthermore, $\Nb^H(x_1)$ must be complete.
  Hence, since $x_0 \tut x_1$ and $x_2 \tut x_2$ in $H$, we must also have $x_0 \tut x_2$ in $H$.
  But then $x_0 \tut x_2 \tut x_3 \dots \tut x_n$ would have been a shorter anterior path from $a$ to $b$ in $H$, a contradiction.

  So the anterior path must have a directed edge.
  Let $k$ be the smallest integer such that $x_k \tuh x_{k+1}$ on $\pi$.
  Note that such $k$ exists and $0 \le k < n$.
  Then the subpath of $\pi$ between $x_k$ and $x_n$ is an anterior path that starts with a directed edge, so there must actually be a directed path from $x_k$ to $x_n = b$ in $H$ by Lemma~\ref{lem:sigmamag_anterior_path_is_ancestral}.
  The (possibly trivial) subpath of $\pi$ between $x_0$ and $x_k$ can only contain undirected edges, by assumption.
  The presence of the edge $b \suh a$ then implies the presence of edges (of the same type) $b \suh x_0$, $b \suh x_1$, \dots, $b \suh x_k$.
  So we would have an edge $b \suh x_k$ and a directed path from $x_k$ to $b$ in $H$, a contradiction.
\end{proof}

Figure~\ref{fig:invalid_sigmamags} provides an example of a mixed graph that satisfies all conditions of a $\sigma$-iMAG except the maximality.
\begin{figure}\centering
  \begin{tikzpicture}
    \begin{scope}
%      \node at (-2,0) {(d)};
      \node[ndout] (X1) at (-1,0) {$i$};
      \node[ndout] (X2) at (-1,2) {$j$};
      \node[ndout] (X3) at (1,2) {$k$};
      \node[ndout] (X4) at (1,0) {$l$};
      \draw[huh] (X1) edge (X2);
      \draw[huh] (X2) edge (X3);
      \draw[huh] (X3) edge (X4);
      \draw[tuh] (X2) edge (X4);
      \draw[tuh] (X3) edge (X1);
    \end{scope}
  \end{tikzpicture}
  \caption{This mixed graph is not a valid $\sigma$-iMAG, because it is not maximal: it has an inducing path $i \huh j \huh  k \huh l$ while $i$ and $l$ are non-adjacent.\label{fig:invalid_sigmamags}}
\end{figure}

The property is an elementary consequences of the definition of $\sigma$-iMAG that we will make use of frequently.
\begin{Lem}\label{lem:sigmamag_anterior_path_is_ancestral}
  Let $H$ be a $\sigma$-iMAG with input nodes $V$ and output nodes $J$.
  Let $a, b \in V \cup J$ be two nodes in $H$.
  If there exists an anterior path from $a$ to $b$ in $H$ that starts with a directed edge, then $a \in \Anc^H(b)$.
\end{Lem}
\begin{proof}
  There exist anterior paths from $a$ to $b$ that start with a directed edge.
  Let $x_0 \tuh x_1 \tus \dots \tus x_n$ be such a path of minimal length.
  If $n=1$ we are done.
  If $n>1$, then consider the edge $x_1 \tus x_2$. 
  If it were undirected, then we would have $x_0 \tuh x_1 \tut x_2$ and hence $x_0 \tuh x_2$; in that case $x_0 \tuh x_2 \tus \dots \tus x_n$ would be a shorter anterior path from $a$ to $b$ that starts with a directed edge, a contradiction.
  Hence it must be directed.
  We can repeat this reasoning and conclude that all edges must be directed.
  Hence $x_0 \in \Anc^H(b)$, that is, $a \in \Anc^H(b)$.
\end{proof}

\subsection{$\sigma$-iMAGs represent CDMGs}

The purpose of $\sigma$-iMAGs is to represent a set of CDMGs. 
That can be seen as a special case of the more general Definition~\ref{def:mixed_graph_represents_cdmg}.
Also, the more general Lemma~\ref{lem:mixed_graph_represents_CDMG} shows some elementary properties of a CDMG that are encoded in a mixed graph that represents it.

The definition of $\sigma$-iMAGs was designed to be a characterization of the mixed graphs (with certain edge types) that represent some CDMG.
\begin{Prp}\label{prp:sigmamag_represents_cdmg_is_valid}
  Let $H = \lt J,V,E\rt$ be a mixed graph with input nodes $J$, output nodes $V$ and edges $E$ of the types $\{\tuh,\hut,\huh,\tut\}$.
  Let $G$ be a CDMG with input nodes $J$ and output nodes $V^+ = V \dcup S$.
  If $H$ represents $G$ given $S$, then $H$ is a $\sigma$-iMAG.
\end{Prp}
\begin{proof}
  By assumption: there is at most one edge between two distinct nodes; there are no edges between a node and itself;
  no input node is adjacent to any other input node. 
  %if there is an edge $j \sus v$ in $H$ with $j \in J, v \in V$ then it must be of the form $j \tus v$.
  
  We show that $H$ is $\sigma$-ancestral. 
  First, suppose that $H$ contained a directed path $i = v_0 \tuh \dots \tuh v_n = j$.
  By Lemma~\ref{lem:mixed_graph_represents_CDMG}, then $i \in \Anc_G(\{j\} \cup S)$.
  If $H$ contained an edge $j \suh i$, this would imply $i \notin \Anc_G(\{j\} \cup S)$, a contradiction. 
  Hence such edges cannot occur.
  This means that no directed cycles and no almost directed cycles occur in $H$.
  %The remaining statement to prove that $H$ is $\sigma$-ancestral follows by Lemma~\ref{lem:into_undirected_edge}.
%  Second, suppose that an unshielded triple of the form $i \suh j \tut k$ occured in $H$.
%  Since $j \in \Anc_G(\{k\} \cup S)$ but $j \notin \Anc_G(\{i\} \cup S)$, we must have $j \in \Anc_G(k)$.
%  Also, $k \in \Anc_G(\{j\} \cup S)$.
%  Since $j \in \Anc_G(k)$ and $j \notin \Anc_G(S)$, we must have $k \in \Anc_G(j)$.
%  But then $j \in \Sc_G(k)$, which is not possible according to Lemma~\ref{lem:induced_path_into_scc}.
%  Hence such unshielded triples cannot occur in $H$.
%  Together, this shows that $H$ is $\sigma$-ancestral.
  It remains to show that no unshielded triple of the form $i \suh j \tut k$ can occur in $H$.
  We prove this by contradiction.
  Since $j \in \Anc_G(\{k\} \cup S)$ but $j \notin \Anc_G(\{i\} \cup S)$, we must have $j \in \Anc_G(k)$.
  Also, $k \in \Anc_G(\{j\} \cup S)$.
  Since $j \in \Anc_G(k)$ and $j \notin \Anc_G(S)$, we must have $k \in \Anc_G(j)$.
  But then $j \in \Sc_G(k)$, which gives a contradiction with Lemma~\ref{lem:inducing_path_into_scc}.
  Hence there must be an edge $i \sus k$ in $H$.
  Now if $i \tuh j$, then $i \in \Anc_G(j)$ or $i \in \Anc_G(S)$, hence $i \in \Anc_G(k)$ or $i \in \Anc_G(S)$, hence the edge between $i$ and $k$ must be of the form $i \tuh k$.
  Otherwise, if $i \huh k$, then $i \notin \Anc_G(j)$ and $i \notin \Anc_G(S)$, hence $i \notin \Anc_G(k)$ and $i \notin \Anc_G(S)$, hence the edge between $i$ and $k$ must be of the form $i \huh k$.
  Finally, if $l_1 \tut j$ and $l_2 \tut j$ with $l_1 \ne l_2$ then both $l_1, l_2 \in \Sc_G(j)$ (since $j \notin \Anc_G(S)$).
  Hence, $l_1 \tut l_2$ must be in $H$.
  Similarly, if $l_1 \tut k$ and $l_2 \tut k$ with $l_1 \ne l_2$ then both $l_1, l_2 \in \Sc_G(j)$, and hence $l_1 \tut l_2$ must be in $H$.

  We continue to show that $H$ is maximal.
  Suppose there is a inducing path $\mu$ in $H$ between two distinct nodes $u,v \in V \cup J$. 
  We first show that this implies that $\{u,v\} \not\subseteq J$.
  If $u,v \in J$, then a inducing path between them cannot consist of a single edge, because all node pairs in $J$ are non-adjacent in $H$ by assumption.
  Hence it must be of the form $u \suh w_1 \dots w_n \hus v$ (with $n \ge 1$) with $w_1, w_n \in V$. 
  But since all nodes $w_1,\dots,w_n$ are colliders, none of them is in $\Anc_G(S)$.
  Since all nodes $w_1,\dots,w_n$ are in $\Anc_H(\{u,v\})$, by Lemma~\ref{lem:mixed_graph_represents_CDMG}, they must all be in $\Anc_G(\{u,v\})$.
  Since $G$ is a CDMG, they must all be in $\{u,v\} \subseteq J$, a contradiction.
  %Both $w_1$ and $w_n$ must be in $\Anc_H(\{u,v\})$, and hence $w_1, w_n \in \Anc_G(\{u,v\} \cup S)$ by Lemma~\ref{lem:mixed_graph_represents_CDMG}.
  %Both $w_1$ and $w_n$ are colliders, hence $w_1 \notin \Anc_G(\{u\} \cup S)$ and $w_n \notin \Anc_G(\{v\} \cup S)$.
  %Therefore, $w_1 \in \Anc_G(v)$ and $w_n \in \Anc_G(w)$.
  %This gives a contradiction with Lemma~\ref{lem:inducing_path_out_of_input}.
  Hence, there cannot be a inducing path in $H$ between two nodes in $J$.

  Every edge $i \sus j$ on $\mu$ corresponds with a $\sigma$-inducing path $\pi_{ij}$ given $S$ in $G$ between $i$ and $j$.
  By Lemma~\ref{lem:mixed_graph_represents_CDMG}, these $\sigma$-inducing paths can be chosen to be into $i$ if the edge is $i \hus j$, into $j$ if the edge is $i \suh j$, and both into $i$ and $j$ if the edge is $i \huh j$. 
  Concatenate all $\pi_{ij}$ (following the edge ordering of $\mu$) into a walk $\pi$ in $G$ between $u$ and $v$.
  Every non-endpoint node on $\mu$ is a collider on $\mu$, by assumption.
  By construction, these nodes then become colliders on $\pi$.
  Since colliders on $\mu$ are in $\Anc_H(\{u,v\})$ by assumption, they are in $\Anc_G(\{u,v\} \cup S)$ by Lemma~\ref{lem:mixed_graph_represents_CDMG}.
  All colliders on some $\pi_{ij}$ are in $\Anc_G(\{i,j\}\cup S)$.
  Hence, they are in $\Anc_G(\{u,v\}\cup S)$.
  So, all colliders on $\pi$ are in $\Anc_G(\{u,v\}\cup S)$.
  All non-endpoint non-colliders on some $\pi_{ij}$ are unblockable, and therefore all non-endpoint non-colliders on $\pi$ are unblockable.
  Therefore, $\pi$ is a $\sigma$-inducing walk given $S$ in $G$.
  So there must also be a $\sigma$-inducing path given $S$ in $G$ between $u$ and $v$.
  Because $H$ represents $G$ given $S$, we conclude that $u,v$ must be adjacent in $H$.
\end{proof}
\Joris{Note that we can propagate the new condition on $\sigma$-MAGs to $\sigma$-PAGs as well and extend the proof of the analogous proposition for $\sigma$-PAGs.
Then, the proof above becomes just a special case of the more general result.

Here is the novel piece:
\begin{Lem}
  Let $G$ be a CDMG with input nodes $J$ and output nodes $V^+ = V \dcup S$ and $H$ a $\sigma$-MAG that represents $G$ given $S$.
  If $a \tuh b \tut c$ in $H$, then $a \tuh c$ in $H$; if $a \huh b \tut c$ in $H$, then $a \huh c$ in $H$.
\end{Lem}
\begin{proof}
  Suppose $a \suh b \tut c$ in $H$.
  Then there exists a $\sigma$-inducing path between $a$ and $b$ in $G$ given $S$ that is into $b$.
  We can extend it with a directed path from $b$ to $c$ into a $\sigma$-inducing path between $a$ and $c$ in $G$ given $S$.
  Hence, $a$ and $c$ must be adjacent in $H$.
  Furthermore, $b,c$ must lie in the same strongly-connected component of $G$, and $b,c \notin \Anc_G(S)$.
  Since $b \notin \Anc_G(a)$, also $c \notin \Anc_G(a)$. 
  So the edge $a \sus c$ must be of the form $a \suh c$.
  If $a \suh b$ is actually $a \tuh b$ then $a \in \Anc_G(b)$; since also $b \in \Anc_G(c)$, the edge $a \suh c$ must be $a \tuh c$.
  If $a \suh b$ is actually $a \huh b$ then $a \notin \Anc_G(b)$; since $c \in \Anc_G(b)$, the edge $a \suh c$ must be $a \huh c$.
\end{proof}
}

The notion of `anterior' was proposed by \cite{RichardsonSpirtes02} for ancestral graphs, and we (trivially) extended its definition to $\sigma$-MAGs.
%\begin{Def}\label{def:sigmamag_anterior}
%  %Let $G$ be a CDMG with input nodes $J$ and output nodes $V^+ = V \dcup S$ and $H$ a $\sigma$-MAG that represents $G$ given $S$.
%  Let $H=\lt J,V,E \rt$ be a mixed graph with input nodes $J$, output nodes $V$ and edges $E$ of the types $\{\tuh,\hut,\huo,\huh,\ouo,\ouh,\tuo,\out,\tut\}$.
%  If there exists a walk in $H$ of the form $v_0 \tus \cdots \tus v_n$ such that each edge is of the form $v_k \tuh v_{k+1}$ or $v_k \tut v_{k+1}$, we say that \emph{$v_0$ is anterior to $v_n$}.
%  For $b \in V \cup J$ we denote by $\Ant^H(b)$ the set of nodes $a \in V \cup J$ that are anterior to $b$.
%  For $B \subseteq V \cup J$ we define $\Ant^H(B) := \bigcup_{b\in B} \Ant^H(b)$.
%\end{Def}
The basic property of this notion is:
\begin{Lem}\label{lem:sigmamag_anterior}
  Let $G$ be a CDMG with input nodes $J$ and output nodes $V^+ = V \dcup S$ and $H$ a $\sigma$-MAG that represents $G$ given $S$.
  Then for $a \in V \cup J$, $B \subseteq V \cup J$:
  $$a \in \Ant^H(B) \iff a \in \Anc^G(B \cup S).$$
\end{Lem}
\begin{proof}
  Suppose $a \in \Ant^H(b)$ for some $b \in B$.
  Then there exists an anterior path $v_0 \tus v_1 \tus \dots \tus v_n$ in $H$ with $a = v_0$ and $b = v_n$.
  For every edge $v_k \tus v_{k+1}$, we have $v_k \in \Anc^G(\{v_{k+1}\} \cup S)$.
  By transitivity, $v_0 \in \Anc^G(\{v_n\} \cup S)$, that is, $a \in \Anc^G(\{b\} \cup S)$.

  Suppose $a \in \Ant^G(\{b\} \cup S)$ for some $b \in B$.
  Then there exists a directed path $v_0 \tuh v_1 \tuh \dots \tuh v_n$ in $G$ with $a = v_0$ and $b = v_n$.
  Each edge on the path is also present in $H$, and will be of the form $v_k \tus v_{k+1}$.
  Hence $a \in \Ant^H(\{b\})$.
\end{proof}

\Joris{Actually, all the above can be easily pushed towards the $\sigma$-PAGs section and then $\sigma$-iMAGs are just a special case of $\sigma$-PAGs.
The only thing we may want to keep is the following result.}

The following result is due to Binghua Yao [private communication] and myself:\footnote{Binghua realized that $\sigma$-MAGs representing a CDMG
have the additional property `$a \tuh b \tut c$ in $H$ implies $a \tuh c$ in $H$, and $a \huh b \tut c$ in $H$ implies $a \huh c$ in $H$; in both cases, $\Nb^H(b)$ and $\Nb^H(c)$ must be complete'
and that together with the other properties, this characterizes the class of $\sigma$-MAGs representing a CDMG.}
\begin{Prp}\label{lem:representing_CDMG_of_sigmamag}
  Let $H$ be a $\sigma$-iMAG with input nodes $J$ and output nodes $V$.
  Then there exists a CDMG $G$ with input nodes $J$ and output nodes $V^+ = V \dcup S$ such that $H$ represents $G$ given $S$.
\end{Prp}
\begin{proof}
  \Joris{TODO: cleanup. Can we think of a more convenient construction? At least it looks as if this one is fairly 'minimal' and could be extended in various ways.}

  We construct the CDMG $G$ as follows.
  It has nodes input nodes $J$ and output nodes $V^+$ with $V^+ = V \cup S$, where 
  %$S := \{s_{\{x,y\}} : x \tut y \in H\}$.
    $$S := \{s_{\{x,y\}} : x \tut y \in H \land (\Nb^H(x) \text{ is incomplete} \lor \Nb^H(y) \text{ is incomplete}) \}.$$
  We add all edges $x \suh y$ for $x,y \in V\cup J$ that are present in $H$ to $G$.
  For each undirected edge $x \tut y$ in $H$: if $\Nb^H(x)$ is incomplete or if $\Nb^H(y)$ is incomplete, we add $x \tuh s_{\{x,y\}} \hut y$ to $G$;
  if both $\Nb^H(x)$ and $\Nb^H(y)$ are complete, we add both directed edges $x \tuh y$ and $x \hut y$ to $G$.
  Since all directed edges in $H$ are also present in $G$, for $x,y \in V \cup J$, $x \in \Anc_H(y) \implies x \in \Anc_{G}(y)$.
  Further, for $x,y \in V \cup J$, $x \in \Anc_G(y)$ implies $x \in \Ant^G(y)$.
  Indeed, by construction of $G$, a directed edge $a \tuh b$ in $G$ implies that $a \tuh b$ or $a \tut b$ must be in $H$.

  For showing that $H$ represents $G$ given $S$, we need to show:
  \begin{enumerate}
    \item For all $x \in V \cup J, y \in V$: $x \sus y$ is in $H$ if and only if there is a $\sigma$-inducing path between $x$ and $y$ given $S$ in $G$;
    \item If $x \hus y$ in $H$ then $x \notin \Anc_{G}(\{y\} \cup S)$;
    \item If $x \tus y$ in $H$ then $x \in \Anc_{G}(\{y\} \cup S)$;
  \end{enumerate}
  
  \Joris{replace $z$ by $a$ and $x$ by $b$ below}
  We now show: if $z \hus x$ in $H$, then $z \notin \Anc_{G}(S)$.
  Assume $z \hus x$ in $H$.
  Suppose $z \in \Anc_{G}(S)$.
  Then there must be a shortest directed path $z = z_0 \tuh \dots \tuh z_k \tuh s$ in $G$ such that $s \in S$, and all other nodes on the directed path are in $V \cup J$ (possibly $k=0$).
  This must correspond with an anterior path from $z$ to $w$ in $H$.
  As $z_k \tuh s$ in $G$ with $s \in S$, by construction of $G$, there must be a $w$ in $H$ such that $z_k \tut w$ in $H$ and $\Nb^H(z_k)$ is incomplete or $\Nb^H(w)$ is incomplete.
  There cannot be any edge $v \suh z_k$ in $H$, because if it were, then $v \suh z_k \tut w$ would be in $H$, and hence both $\Nb^H(z_k)$ and $\Nb^H(w)$ would be complete, a contradiction.
  Hence, if $k \ge 1$, the edge $z_{k-1} \tus z_k$ in $H$ cannot be $z_{k-1} \tuh z_k$ but must be $z_{k-1} \tut z_k$.
  If $k \ge 2$, the edge $z_{k-2} \tut z_{k-1}$ must be in $H$, because if instead $z_{k-2} \tuh z_{k-1}$ were in $H$, then also $z_{k-2} \tuh z_k$ in $H$ and hence in $G$, a contradiction.
  We can use this reasoning repeatedly to conclude that $z_0 \tut z_1 \tut \dots \tut z_k$ must be in $H$.
  But since we assumed $x \suh z$ in $H$, this then implies $x \suh z_k$ is also in $H$, a contradiction.

  We now show: if $z \hus x$ in $H$, then $z \notin \Anc_G(x)$.
  First assume $z \hut x$ in $H$. 
  Then also $z \hut x$ in $G$. 
  If $z \in \Anc_G(x)$ then there is a directed path from $z$ to $x$ in $G$.
  The directed path in $G$ corresponds to an anterior path in $H$ from $z$ to $x$.
  But then we have in $H$ an anterior path from $z$ to $x$ and an edge $z \hut x$, a contradiction.
  Hence $z \notin \Anc_G(x)$.

  So together we have shown that $z \hus x$ in $H$ implies $z \notin \Anc_G(\{x\} \cup S)$.
  Now suppose $a \tus b$ in $H$. 
  Then either $a \tut b$ in $H$ or $a \tuh b$ in $H$.
  In the first case, by construction of $G$ we conclude $a \in \Anc_G(\{b\} \cup S)$.
  In the second case, by construction of $G$ we have $a \in \Anc_G(b)$.

  Finally, adjacency.
  Let $x \in V \cup J$, $y \in V$.
  Suppose $x$ and $y$ are adjacent in $H$.
  Then they are either adjacent in $G$ (in which case there exists a $\sigma$-inducing path between $x$ and $y$ given $S$ in $G$), or we have $x \tuh s_{\{x,y\}} \hut y$ in $G$ (in which case that would be a $\sigma$-inducing path between $x$ and $y$ given $S$ in $G$).

  Suppose $x$ and $y$ are not adjacent in $H$.
  We must show that in $G$ there is no $\sigma$-inducing path between $x$ and $y$ given $S$.
  Suppose there is one.
  Pick such a path $\pi$ of minimal length, say $x = x_0 \sus \dots \sus x_n = y$.
  It must be a path in which every collider is in $\Anc_{G}(\{x,y\} \cup S)$, and every non-endpoint non-collider is unblockable.

  Suppose $\pi$ contains a node of $S$.
  Then it has to be of the form $x \dots a \tuh s_{\{a,b\}} \hut b \dots y$ with $a \tut b$ in $H$.
  Here, $s_{\{a,b\}}$ lies in another strongly-connected component than $a$, and therefore $a$ is blockable.
  Similarly, $b$ is blockable.
  Since only the endpoint nodes can be blockable non-colliders, the path must be $x \tuh s_{\{x,y\}} \hut y$ with $x \tut y$ in $H$. 
  Contradiction, since $x$ and $y$ were assumed to be non-adjacent in $H$.
  We conclude that the path can only contain nodes in $V \cup J$.
  Also, there must be a corresponding path $\tilde{\pi}$ of the form $x_0 \dots x_n$ in $H$ (with the same nodes but possibly different edges).

  Suppose $\pi$ contains a non-endpoint non-collider.
  If it is $x_{k-1} \sus x_k \tuh x_{k+1}$ then $x_k$ and $x_{k+1}$ must be in the same strongly-connected component of $G$, otherwise $x_k$ would be blockable.
  Hence we must have $x_k \tut x_{k+1}$ in $H$, and both $\Nb^H(x_k)$ and $\Nb^H(x_{k+1})$ complete.

  If we have $x_{k-1} \huh x_k$ on $\pi$, then also $x_{k-1} \huh x_k$ in $H$, and therefore $x_{k-1} \huh x_{k+1}$ in $H$, and hence $x_{k-1} \huh x_{k+1}$ in $G$, and hence there exists a shorter $\sigma$-inducing path between $x$ and $y$ given $S$ in $G$, a contradiction.\\
  If we have $x_{k-1} \tuh x_k$ on $\pi$, then either $x_{k-1} \tuh x_k$ in $H$, or $x_{k-1} \tut x_k$ in $H$ and $x_{k-1}$ and $x_k$ are in the same strongly-connected component of $G$ and both $\Nb^H(x_{k-1})$ nd $\Nb^H(x_k)$ are complete.\\
  If we have $x_{k-1} \hut x_k$ on $\pi$, then either $x_{k-1} \hut x_k$ in $H$, or $x_{k-1} \tut x_k$ in $H$ and $x_{k-1}$ and $x_k$ are in the same strongly-connected component of $G$ and both $\Nb^H(x_{k-1})$ nd $\Nb^H(x_k)$ are complete.\\
  If $x_{k-1} \tuh x_k$ in $H$, then also $x_{k-1} \tuh x_{k+1}$ in $H$, and hence $x_{k-1} \tuh x_{k+1}$ in $G$.
  Since $x_{k-1} \tuh x_k$ on $\pi$, either $x_{k-1}$ is an endpoint or $x_{k-1}$ and $x_k$ must be in the same strongly-connected component of $G$.
  In the first case, there would exist a shorter $\sigma$-inducing path $x_{k-1} \tuh x_{k+1} \dots$ in $G$, a contradiction.
  In the second case, we get a contradiction because it would require $x_{k-1} \tut x_k$ in $H$.\\
  If $x_{k-1} \hut x_k$ in $H$, then also $x_{k-1} \hut x_k$ in $G$, and then we must have $x_{k-1} \hut x_k \tuh x_{k+1}$ on $\pi$.
  But then all three of $x_{k-1}$, $x_k$ and $x_{k+1}$ must be in the same strongly-connected component of $G$, which contradicts the arrowhead $x_{k-1} \hut x_k$ in $H$.
  If $x_{k-1} \tut x_k$ in $H$ and $x_{k-1}$ and $x_k$ are in the same strongly-connected component of $G$ and both $\Nb^H(x_{k-1})$ nd $\Nb^H(x_k)$ are complete, then also $x_{k-1} \tut x_{k+1}$ in $H$ (because $\Nb^H(x_k)$ is complete) and therefore both $x_{k-1} \tuh x_{k+1}$ and $x_{k-1} \hut x_{k+1}$ in $G$ (because both $\Nb^H(x_{k-1})$ and $\Nb^H(x_{k+1})$ are complete); then we could have shortened $\pi$ by replacing $x_{k-1} \tuh x_k \tuh x_{k+1}$ by $x_{k-1} \tuh x_{k+1}$ or $x_{k-1} \hus x_k \tuh x_{k+1}$ by $x_{k-1} \tuh x_{k+1}$
  (note that if this replacement would turn $x_{k-1}$ into a collider, then it would be in $\Anc_G(\{x,y\} \cup S)$, because \emph{any} node on $\pi$ must be in $\Anc_G(\{x,y\} \cup S)$).

  Note that each edge on $\pi$, except possibly for the ones at the end points, must be bidirected, and hence be present in $H$ as well.
  For each undirected edge on $\tilde{\pi}$, its corresponding nodes must have complete neighborhoods in $H$.
  So if $\tilde{\pi}$ would be of the form $x \tut z \tut y$ we conclude that also $x \tut y$ in $H$, a contradiction.
  Also, if $\tilde{\pi}$ were of the form $x \tut z \hus y$ we arrive at $x \hus y$ in $H$, a contradiction.
  Similarly, if $\tilde{\pi}$ were of the form $x \suh z \tut y$ we arrive at $x \suh y$ in $H$, a contradiction.
  Hence if $\tilde{\pi}$ consists of two edges, it must be of the form $x \suh z \hus y$.
  If $\tilde{\pi}$ consists of three edges, then if it is $x \tut a \huh b \tut y$, we must have $x \huh y$ in $H$, a contradiction.
  So if $\tilde{\pi}$ consists of three edges and contains undirected edges, it must be either of the form $x \suh a \huh b \tut y$ or $x \tut a \huh b \hus y$.
  In general, if $\tilde{\pi}$ is of the form $x_0 \tut x_1 \hus x_2 \dots x_n$ then we shorten it to $x_0 \hus x_2 \dots x_n$.
  Then, if (the shortened) $\tilde{\pi}$ is of the form $x_0 \dots x_{n-2} \suh x_{n-1} \tut x_n$ we shorten it to $x_0 \dots x_{n-2} \suh x_n$.
  The (shortened) $\tilde{\pi}$ is then a collider path in $H$.
  Each collider is in $\Anc_G(\{x,y\} \cup S)$, and not in $\Anc_G(S)$ (because of the arrowheads in $H$), hence in $\Anc_G(\{x,y\})$.
  Therefore, it must be in $\Ant^H(\{x,y\})$.
  Now by Lemma~\ref{lem:sigmamag-inducing_path_anterior}, there must exist an inducing path in $H$ between $x$ and $y$.
  Because $H$ is assumed to be maximal, this contradicts the assumed non-adjacency of $x$ and $y$ in $H$.
\end{proof}

\subsection{$im$-separation in $\sigma$-iMAGs}

Perhaps we can prove a more general result that allows us to quickly see this.
\begin{Lem}\label{lem:sigmamag_collider_anterior_or_ancestral}
  Let $H$ be a $\sigma$-iMAG with input nodes $V$ and output nodes $J$.
  Let $x_0 \dots x_n$ be a walk in $H$ that contains a collider $x_k$ that is in $\Ant^H(C)$ for some subset $C \subseteq V \cup J$.
  Then we can replace the subpath $x_{k-1} \suh x_k \hus x_{k+1}$ by $x_{k-1} \suh \tilde{x}_k \hus x_{k+1}$, where $\tilde{x}_k \in \Anc^H(C)$ and the edge marks at $x_{k-1}$ and at $x_{k+1}$ are preserved.
%  If $x_k \in \Anc^H(C)$ then obviously $x_k \in \Ant^H(C)$ without the need for replacement.
\end{Lem}
\begin{proof}
  There exists a shortest anterior path from $v_k$ to $C$.
  Suppose $w_0 \tus \dots \tus w_r$ with $w_0 = v_k$ and $w_r \in C$ is such a path.
  If $v_k \in C$ then there is nothing left to do because then also $v_k \in \Anc^H(C)$.
  If $v_k \notin C$, then the anterior path consists of at least one edge.
  Suppose the anterior path starts with an undirected edge: $w_0 \tut w_1 \tus \dots \tus w_r$ (with possibly $r=1$).
  Since $v_{k-1} \suh v_k \tut w_1$ in $H$, we have $v_{k-1} \suh w_1$ in $H$ with $v_{k-1} \suh w_1$ the same edge type as $v_{k-1} \suh v_k$.
  Similarly, since $v_{k+1} \suh v_k \tut w_1$ in $H$, we have $v_{k+1} \suh w_1$ in $H$ with $v_{k+1} \suh w_1$ the same edge type as $v_{k+1} \suh w_1$.
  We can now replace the subpath $v_{k-1} \suh v_k \hus v_{k+1}$ of $\pi$ with the subpath $v_{k-1} \suh w_1 \hus v_{k+1}$.
  Since the only change is the node $w_1$ instead of $v_k$, while all edges remain the same, and $w_1 \in \Ant^H(C)$, this new path has the same property of $\pi$.
  We can extend this reasoning actually to the first undirected segment rather than just the first undirected edge of the anterior path $w_0 \tus \dots \tus w_r$.
  We set $\tilde{v}_k := w_s$.
  
  We have an anterior path $w_s \tus \dots \tus w_r$ that starts with a directed edge in case $s < r$.
  By Lemma~\ref{lem:sigmamag_anterior_path_is_ancestral} we conclude that $\tilde{v}_k = w_s \in \Anc^H(w_r) \subseteq \Anc^H(C)$.
  If $s=r$ it is obvious that $\tilde{v}_k = w_s = w_r \in \Anc^H(C)$.
\end{proof}

The cases in Figure~\ref{fig:sigmapags_funky_separations} are helpful for the following proof.
\begin{figure}\centering
  \begin{tikzpicture}
    \begin{scope}
      \node at (-2,1.5) {(a)};
      \node[ndout] (X1) at (-1,1.5) {$i$};
      \node[ndout] (X2) at (-1,0) {$j$};
      \node[ndout] (X3) at (1,1.5) {$k$};
      \node[ndout] (X4) at (1,0) {$l$};
      \draw[tuh] (X1) edge (X2);
      \draw[tuh] (X3) edge (X4);
      \draw[tuh,bend left] (X2) edge (X4);
      \draw[tuh,bend left] (X4) edge (X2);
    \end{scope}
    \begin{scope}[xshift=5cm]
      \node at (-2,1.5) {(b)};
      \node[ndout] (X1) at (-1,1.5) {$i$};
      \node[ndout] (X2) at (-1,0) {$j$};
      \node[ndout] (X3) at (1,1.5) {$k$};
      \node[ndout] (X4) at (1,0) {$l$};
      \draw[tuh] (X1) edge (X2);
      \draw[tuh] (X3) edge (X2);
      \draw[tuh,bend left] (X2) edge (X4);
      \draw[tuh,bend left] (X4) edge (X2);
    \end{scope}
    \begin{scope}[xshift=11cm]
      \node at (-2,1.5) {(c)};
      \node[ndout] (X1) at (-1,1.5) {$i$};
      \node[ndout] (X2) at (-1,0) {$j$};
      \node[ndout] (X3) at (1,1.5) {$k$};
      \node[ndout] (X4) at (1,0) {$l$};
      \draw[tuh] (X1) edge (X2);
      \draw[tuh] (X3) edge (X4);
      \draw[tut] (X2) edge (X4);
      \draw[tuh] (X1) edge (X4);
      \draw[tuh] (X3) edge (X2);
    \end{scope}
  \end{tikzpicture}
  \caption{Useful examples for relating the existence of $\sigma$-open paths in a CDMG with the existence of $m$-open paths in a $\sigma$-MAG that represents it.
  (a) and (b) are two CDMGs that are both represented by the $\sigma$-MAG in (c).
  In (a), $i \isPerp k$. 
  In (b), a $\sigma$-open path $i \tuh j \hut k$ given $\{l\}$ exists.
  In (c), there is a $d$-open path $i \tuh j \tut l \hut k$, but no $m$-open path between $i$ and $k$.
  Furthermore, in (c), the path $i \tuh j \hut k$ is $m$-open given $\{l\}$ because we use $\Ant$ instead of $\Anc$.
  \label{fig:sigmapags_funky_separations}}
\end{figure}

We extend the notion of $m$-blocking defined for MAGs \cite{RichardsonSpirtes02} to $\sigma$-MAGs:
\begin{Def}\label{def:sigmamag_blocking}
  Let $G$ be a CDMG with input nodes $J$ and output nodes $V^+ = V \dcup S$ and $H$ a $\sigma$-MAG that represents $G$ given $S$.
  Let $\pi$ be a walk in $H$:
  \[ \pi =\lp  v_0 \sus \cdots \sus v_n \rp.  \]
  Let $C \subseteq J \cup V$. We say that the walk $\pi$ is:
  \begin{enumerate}
    \item \emph{$C$-$m$-open} (or \emph{$m$-open given $C$}) if and only if:
      \begin{enumerate}[label=\roman*)]
        %\item all colliders $v_k$ on $\pi$ are in $\Anc^H(C)$, \emph{and}
        \item all colliders $v_k$ on $\pi$ are in $\Ant^H(C)$, \emph{and}
        \item all non-colliders $v_k$ on $\pi$ are \emph{not} in $C$, \emph{and}:
        \item it contains no subwalks of the form $v_{k-1} \suh v_k \tut v_{k+1}$ or $v_{k-1} \tut v_k \hus v_{k+1}$.
      \end{enumerate}
    \item \emph{$C$-$m$-blocked} (or \emph{$m$-blocked given $C$}) if and only if:
      \begin{enumerate}[label=\roman*)]
        %\item there exists a collider $v_k$ on $\pi$ that is \emph{not} in $\Anc^H(C)$
        \item there exists a collider $v_k$ on $\pi$ that is \emph{not} in $\Ant^H(C)$, \emph{or}:
        \item there exists a non-collider $v_k$ on $\pi$ in $C$, \emph{or}:
        \item there exists a subwalk of the form $v_{k-1} \suh v_k \tut v_{k+1}$ or $v_{k-1} \tut v_k \hus v_{k+1}$.
      \end{enumerate}  
  \end{enumerate}
  Hence the walk is $C$-$m$-open if and only if it is not $C$-$m$-blocked.
\end{Def}

Lemma~\ref{lem:ancestral_relations_open_path} can be generalized to $\sigma$-MAGs:
\begin{Lem}\label{lem:sigmamag_ancestral_relations_open_path}
  Let $G$ be a CDMG with input nodes $J$ and output nodes $V^+ = V \dcup S$ and $H$ a $\sigma$-MAG that represents $G$ given $S$.
  Let $i, j \in J \cup V$ and $Z \subseteq J \cup V$.
  If $\pi$ is an $m$-open walk given $Z$ between $i$ and $j$ in $H$, then 
  %every node on $\pi$ is in $\Anc_G(\{i,j\} \cup Z \cup S)$.
%  \Joris{In case you expected the statement ``If $\pi$ is a $Z$-open walk between $i$ and $j$ in $H$, then every node on $\pi$ is in $\Anc_H(\{i,j\} \cup Z)$.'' then this is not true. Indeed, consider $i \tuh s_1 \hut k \tuh s_2 \hut j$. We could however conclude that every node on $\pi$ is in $\Ant^H(\{i,j\} \cup Z)$?}
  every node on $\pi$ is in $\Ant^H(\{i,j\} \cup Z)$.
\end{Lem}
\begin{proof}
  Note that all colliders on $\pi$ are in $\Ant^H(Z)$.
% Note that a collider on $\pi$ cannot be in $\Ant^H(J)$. Because then there must be an anterior path with an incoming arrowhead $\suh v_0 \tus \dots \tus v_n$ with $v_n \in J$, a contradiction.
  Consider now a non-collider $k$ on $\pi$. 
  Then there is an anterior subwalk from $k$ to an endnode of $\pi$, in which case $k \in \Ant_H(\{i,j\})$,
% \Joris{which might be in $J$ in case the anterior path has no incoming arrowhead?}
  or there is an anterior subwalk from $k$ to a collider on $\pi$, in which case $k \in \Ant_H(Z)$,
  So each node on $\pi$ is in $\Ant^H(\{i,j\}\cup Z)$.
\end{proof}
We can also straightforwardly generalize Lemma 17 in \cite{SpirtesRichardson97} to $\sigma$-MAGs.
\begin{Lem}\label{lem:sigmamag_open_path_implies_cdmg_sigma-open_path}
  Let $G$ be a CDMG with input nodes $J$ and output nodes $V^+ = V \dcup S$ and $H$ a $\sigma$-MAG that represents $G$ given $S$.
  Let $Z \subseteq J \cup V$ and
    \[\pi = (q_0, a_1, q_1, \dots q_{n-1}, a_n, q_n)\] 
  be an $m$-open path given $Z$ in $H$.
  Then there exists a $\sigma$-open path given $(Z \cup S)$ in $G$ between $q_0$ and $q_n$.
\end{Lem}
\begin{proof}
  For $i=1,\dots,n$, let $\mu_i$ be a $\sigma$-inducing path in $G$ between $q_{i-1}$ and $q_i$ that is into $q_{i-1}$ if the edge $q_{i-1} \hus q_i$ on $\pi$, and into $q_{i}$ if the edge $q_{i-1} \suh q_i$ on $\pi$ (see Lemma~\ref{lem:mixed_graph_represents_CDMG}).
  These can be concatenated into a walk $\mu = (\mu_1,\dots,\mu_n)$ in $G$ between $q_0$ and $q_n$:
  \[\overunderbraces{&\br{2}{\mu_1}&&&&\br{2}{\mu_{n}}}{&q_0 \cdots \sus & q_1 & \sus \dots \sus q_2 & \sus \cdots \cdots \sus & q_{n-2} \sus \dots \sus & q_{n-1} & \sus \cdots q_n}{&&\br{2}{\mu_2}&&\br{2}{\mu_{n-1}}}\]

  Consider all walks in $G$ between $q_0$ and $q_n$ with the property that 
  \begin{enumerate}
    \item all colliders on it are in $\Anc_{G}(\{q_0,q_n\} \cup Z \cup S)$, and 
    \item each non-endpoint non-collider on it is not in $\{q_0,q_n\} \cup Z \cup S$ or is unblockable.
  \end{enumerate}
  Such walks exist, since the concatenation $\mu = (\mu_1, \dots, \mu_n)$ is one, as we will now show. 

  To show the first property, note that by Lemma~\ref{lem:sigmamag_ancestral_relations_open_path}, all $q_i \in \Anc_G(\{q_0,q_n\} \cup Z \cup S)$.
  This holds in particular for all $q_i$ that are a collider on $\mu$.
  Every internal collider on some $\mu_i$ is in $\Anc_G(\{q_{i-1},q_i\} \cup S)$, and hence also these internal colliders are in $\Anc_G(\{q_0,q_n\} \cup Z \cup S)$.

  For the second property, note that all internal non-colliders on some $\mu_i$ are unblockable by assumption.
  Suppose that $q_i$ ($0 < i < n$) is a non-collider on $\mu$. 
  Then $q_i \notin \{q_0,q_n\}$ because all nodes $q_0,\dots,q_n$ are distinct.
  $q_i$ cannot be a collider on $\pi$, as it would then also be a collider on $\mu$.
  Hence $q_i$ is a non-collider on $\pi$, and therefore not in $Z$.
  Obviously, $q_i \notin S$.
  Hence, $q_i \notin \{q_0,q_n\} \cup Z \cup S$.

  Now let $\nu$ be such a walk with a minimal number of colliders and such that $q_0$ and $q_n$ only occur once on $\nu$.
  We show that all colliders on $\nu$ must be in $\Anc_{G}(Z \cup S)$. 

  Suppose on the contrary the existence of a collider $k$ on $\nu$ that is not in $\Anc^G(Z \cup S)$.
  It is in $\Anc^G(\{q_0,q_n\})$, by assumption.
  Then there exists a directed path in $G$ from $k$ to $q_0$ that does not pass through $q_n$, or there exists a directed path in $G$ from $k$ to $q_n$ that does not pass through $q_0$.
  Without loss of generality, assume the former:
  there is a directed path from $k$ to $q_0$ in $G$ that does not pass through $\{q_n\} \cup Z \cup S$.
  Let $\tau$ be a shortest path (possibly trivial) of that type.
  The directed path $\tau$ from $k$ to $q_0$ can be concatenated with the subwalk of $\nu$ between $k$ and $q_n$ (which is into $k$) into a walk between $q_0$ and $q_n$.
  This walk has the properties above, but with fewer colliders than $\nu$: a contradiction.

%  Therefore, $\nu$ is $\sigma$-open given $Z$ and all occurrences of $q_0$ and $q_n$ as non-endpoint non-collider are unblockable.
  Therefore, $\nu$ is a $\sigma$-open walk given $Z \cup S$ in $G$ between $q_0$ and $q_n$.
  This means that there must also exist a $\sigma$-open path in $G$ between $q_0$ and $q_n$ given $Z \cup S$.
\end{proof}

It is somewhat more involved to generalize its converse, Lemma 18 in \cite{SpirtesRichardson97}, to our more general setting:
\begin{Lem}\label{lem:cdmg_sigma-open_path_implies_sigmamag_open_path}
  Let $G$ be a CDMG with input nodes $J$ and output nodes $V^+ = V \dcup S$ and $H$ a $\sigma$-MAG that represents $G$ given $S$.
  Let $Z \subseteq (J \cup V)$ and $a,b \in J \cup V$.
  If there exists a $\sigma$-open walk given $Z \cup S$ between $a$ and $b$ in $G$, then there exists an $m$-open path given $Z$ between $a$ and $b$ in $H$.
\end{Lem}
\begin{proof}
  Assume (without loss of generality) that $\mu = q_0 \sus q_1 \sus \dots \sus q_n$ is a $\sigma$-open path between $q_0=a$ and $q_n=b$ given $Z \cup S$ in $G$ such that if $q_i \in \Sc^G(q_j)$ for $i < j$, then the subpath of $\mu$ between $q_i$ and $q_j$ is a directed path consisting only of nodes in $\Sc^G(q_i)$.
  For $0 \le i < j \le n$, denote by $\mu(q_i,q_j)$ the subpath of $\mu$ between $q_i$ and $q_j$.
  Note: all blockable non-colliders on $\mu$ are not in $Z \cup S$, all colliders on $\mu$ are in $\Anc_G(Z \cup S)$.

  Consider the subsequence $v^0_{0}, \dots, v^0_{m_0}$ of nodes on $\mu$ that are blockable non-colliders on $\mu$ or colliders that are not in $\Anc_G(S)$, with $v^0_{0} = q_0$ and $v^0_{m_0} = q_n$.
  That is, we skip the unblockable non-colliders and the colliders that are in $\Anc_G(S)$; in other words, we skip all non-endpoint nodes on $\mu$ that cannot be $\sigma$-blocked no matter what $Z$ would be.

  Note that for $0 \le i < m_0$, $\mu(v^0_{i},v^0_{i+1})$ is a $\sigma$-inducing path given $S$.
  Indeed, all non-endpoint non-colliders on it are unblockable, and all colliders are in $\Anc_G(S)$.
  If $\mu(v^0_{i},v^0_{i+1})$ is out of $v^0_{i}$, then $v^0_{i} \in \Anc^G(\{v^0_{i+1}\} \cup S)$.
  Indeed, any non-endpoint nodes on this subpath must be unblockable non-colliders or colliders in $\Anc_G(S)$.
  Hence there exists a subpath $\mu(v^0_{i},q_l)$ such that $q_l$ is a collider in $\Anc_G(S)$ or $q_l = v^0_{i+1}$, and with all non-endpoint nodes being unblockable non-colliders and hence in the same strongly-connected component of $G$ as $v^0_{i}$.
  Similarly, if $\mu(v^0_{i},v^0_{i+1})$ is out of $v^0_{i+1}$, then $v^0_{i+1} \in \Anc^G(\{v^0_{i}\} \cup S)$.

  So there is a path $\pi^0 := v^0_{0} \sus v^0_{1} \sus \dots \sus v^0_{m_0}$ in $H$.
  If $v^0_{i}$ is a non-endpoint non-collider on $\mu$, then it is in $\Anc_G(\{v^0_{i-1}\} \cup S)$ or in $\Anc_G(\{v^0_{i+1}\} \cup S)$ as we saw before, and hence a non-collider on $\pi^0$ as well.
  If $v^0_{i}$ is a collider on $\mu$, then it is not in $\Anc_G(S)$ but it could be $\Anc_G(v^0_{i-1})$ or in $\Anc_G(v^0_{i+1})$.
  The following algorithm removes such colliders.
  \begin{algorithmic}[1]
    \State $k \leftarrow 0$
    \Repeat
      \State $I \leftarrow \{0 < i < m_k : v^{k}_{i} \in \Anc_G(\{v^k_{i-1},v^k_{i+1}\}) \text{ and } \text{$v^k_{i}$ is a collider on $\mu$} \}$
      \If{$I \ne \emptyset$}
        \State form sequence $v^{k+1}$ from $v^k$ by removing some $v^k_{i}$ with $i \in I$
        \State $k \leftarrow k+1$
      \EndIf
      \Until{$I = \emptyset$}
    \State $K \leftarrow k$
  \end{algorithmic}
  Note that in each iteration $k$, $\mu(v^{k+1}_{i-1},v^{k+1}_{i})$ is a $\sigma$-inducing path between $v^{k+1}_{i-1}$ and $v^{k+1}_{i}$ given $S$ in $G$.
  Indeed, concatenating the two paths $\mu(v^k_{i-1},v^k_i)$, $\mu(v^k_i,v^k_{i+1})$ at $v^k_i$, that are both $\sigma$-inducing given $S$, into $v^k_i$, and where $v^k_i \in \Anc_G(\{v^k_{i-1},v^k_{i+1})$, gives the concatenated path $\mu(v^k_{i-1},v^k_{i+1})$ which must be $\sigma$-inducing given $S$.
  Therefore, there is a path $\pi^{k+1} := v^{k+1}_{1} \sus v^{k+1}_{2} \sus \dots \sus v^{k+1}_{m_{k+1}}$ in $H$ (with $m_{k+1} := m_k - 1$).
  Note that we never remove an end-point node, so $v^k_{0} = q_0$ for all $k$, and $v^k_{m_k} = q_n$ for all $k$.
  We will show that the path $\pi^K$ in $H$ is $m$-open given $Z$.

  In the remainder of the proof, we will use the following terminology: we call a node $q_i$ ($0 \le i < n$) on $\mu$ \emph{acyclically right-out-of} if $\mu(q_i,q_{i+1}) = q_i \tuh q_{i+1}$ with $q_{i+1} \notin \Sc_G(q_i)$, and we call a node $q_i$ ($0 < i \le n$) \emph{acyclically left-out-of} if $\mu(q_{i-1},q_i) = q_{i-1} \hut q_i$ with $q_{i-1} \notin \Sc^G(q_i)$.
%  We will prove by induction that for $k=0,\dots,K$: for any $0 \le i < m_k$, if $\mu(v^k_{i},v^k_{i+1})$ is out of $v^k_{i}$ then $v^k_{i} \in \Anc_G(\{v^k_{i+1}\} \cup S)$, and if $\mu(v^k_{i},v^k_{i+1})$ is out of $v^k_{i+1}$ then $v^k_{i+1} \in \Anc_G(\{v^k_{i}\} \cup S)$.
  We will prove by induction that for $0 \le k \le K$: if $v^k_i$ ($0 \le i < m_k$) is acyclically right-out-of then $v^k_{i} \in \Anc_G(\{v^k_{i+1}\} \cup S)$, and if $v^k_i$ ($0 < i \le m_k$) is acyclically left-out-of then $v^k_{i} \in \Anc_G(\{v^k_{i-1}\} \cup S)$.
  The base case $k=0$ has been shown already.
  So suppose the claim holds for $0\le k<K$.
  Suppose $v^{k+1}$ is formed by removing $v^k_{i}$:
  $$\begin{array}{lllllllllll}
    v^k =     & v^k_{0},     & \dots, & v^k_{i-2},     & v^k_{i-1},     & v^k_{i}, & v^k_{i+1},   & v^k_{i+2},     & \dots, & v^k_{m_k} \\ 
    \\
    v^{k+1} = & v^{k+1}_{0}, & \dots, & v^{k+1}_{i-2}, & v^{k+1}_{i-1}, & \cancel{v^k_{i},}         & v^{k+1}_{i}, & v^{k+1}_{i+1}, & \dots, & v^{k+1}_{m_{k+1}}
  \end{array}$$
%  We have to check whether the claim still holds for $v^{k+1}_{i-1}$ (who got a new right neighbor $v^{k+1}_{i} = v^k_{i+1}$) and for $v^{k+1}_{i}$ (who got a new left neighbor $v^{k+1}_{i-1} = v^k_{i-1}$).
%  All other nodes on $v^{k+1}$ keep their neighbors and hence the property is preserved for those nodes.
  We could only have removed $v^k_{i}$ if $v^k_{i} \in \Anc_G(\{v^k_{i-1},v^k_{i+1}\})$.
  Let us first consider the case that $v^k_{i} \in \Anc_G(v^k_{i+1})$.
  Note that it suffices to check whether the claim still holds for the nodes that got a new neighbor in the sequence: $v^{k+1}_{i-1}$ and $v^{k+1}_{i}$.
  \begin{itemize}
    \item Suppose $v^{k+1}_{i-1}$ is acyclically right-out-of on $\mu$. 
      In other words, $v^k_{i-1}$ is acyclically right-out-of on $\mu$.
      By the induction hypothesis, $v^k_{i-1} \in \Anc_G(\{v^k_{i}\} \cup S)$.
      Since we assumed $v^k_{i} \in \Anc_G(v^k_{i+1})$, we conclude $v^k_{i-1} \in \Anc_G(\{v^k_{i+1}\} \cup S)$.
      Hence $v^{k+1}_{i-1} \in \Anc_G(\{v^{k+1}_{i}\} \cup S)$.
      If $v^{k+1}_{i-1}$ is acyclically left-out-of on $\mu$, the claim follows immediately from the induction hypothesis.
%      \begin{itemize}
%        \item Suppose $\mu(v^{k+1}_{i-1},v^{k+1}_{i})$ is seriously out of $v^{k+1}_{i-1}$.
%          Then $\mu(v^k_{i-1},v^k_{i})$ is seriously out of $v^k_{i-1}$.
%          By the induction hypothesis, $v^k_{i-1} \in \Anc_G(\{v^k_{i}\} \cup S)$.
%          Since we assumed $v^k_{i} \in \Anc_G(v^k_{i+1})$, we conclude $v^k_{i-1} \in \Anc_G(\{v^k_{i+1}\} \cup S)$.
%          Hence $v^{k+1}_{i-1} \in \Anc_G(\{v^{k+1}_{i}\} \cup S)$.
%        \item (nothing to show!) Suppose $\mu(v^{k+1}_{i-2},v^{k+1}_{i-1})$ is seriously out of $v^{k+1}_{i-1}$.
%          Then $\mu(v^k_{i-2},v^k_{i-1})$ is seriously out of $v^k_{i-1}$.
%          By the induction hypothesis, $v^k_{i-1} \in \Anc_G(\{v^k_{i-2}\} \cup S)$.
%          Hence $v^{k+1}_{i-1} \in \Anc_G(\{v^{k+1}_{i-2}\} \cup S)$.
%      \end{itemize}
    \item If $v^{k+1}_{i}$ is acyclically right-out-of on $\mu$, the claim follows immediately from the induction hypothesis.
      Suppose $v^{k+1}_{i}$ is acyclically left-out-of on $\mu$.
      In other words, $v^k_{i+1}$ is acyclically left-out-of on $\mu$.
      By the induction hypothesis, $v^{k}_{i+1} \in \Anc_G(\{v^k_{i}\} \cup S)$.
      Remember that we assumed that $v^k_{i} \in \Anc_G(v^k_{i+1})$. 
      Hence, $v^{k}_{i+1} \in \Anc_G(v^k_i)$ would imply $v^k_i \in \Sc_G(v^k_{i+1})$.
      By assumption on the form of $\mu$, this would contradict that $v^k_{i+1}$ is acyclically left-out-of on $\mu$.
      Hence $v^{k}_{i+1} \in \Anc_G(S)$.
      But then also $v^k_i \in \Anc_G(S)$, a contradiction (since all colliders on $\mu$ in $\Anc_G(S)$ were already removed in an earlier stage). %, implying $v^{k+1}_{i} \in \Anc_G(\{v^{k+1}_{i-1}\} \cup S)$.
%      \begin{itemize}
%        \item (nothing to show!) Suppose $\mu(v^{k+1}_{i},v^{k+1}_{i+1})$ is seriously out of $v^{k+1}_{i}$.
%          Then $\mu(v^k_{i+1},v^k_{i+2})$ is seriously out of $v^k_{i+1}$.
%          By the induction hypothesis, $v^k_{i+1} \in \Anc_G(\{v^k_{i+2}\} \cup S)$.
%          Hence $v^{k+1}_{i} \in \Anc_G(\{v^{k+1}_{i+1}\} \cup S)$.
%        \item Suppose $\mu(v^{k+1}_{i-1},v^{k+1}_{i})$ is seriously out of $v^{k+1}_{i}$.
%          That is, $\mu$ is of the form $v^{k+1}_{i-1} \dots v \hut v^{k+1}_{i} \dots$ with $v \notin \Sc_G(v^{k+1}_{i})$ ($v = v^{k+1}_{i-1}$ cannot occur since $v^k_{i}$ lies in between).
%          By the induction hypothesis, $v^{k}_{i+1} \in \Anc_G(\{v^k_{i}\} \cup S)$.
%          If $v^k_{i+1} \notin \Anc_G(S)$ then $v^{k}_{i} \in \Sc_G(v^k_{i+1})$.
%          But then also $v \in \Sc_G(v^k_{i+1})$ by assumption on $\mu$.
%          Contradiction.
%          Hence $v^k_{i+1} \in \Anc_G(S)$.
%          Then certainly $v^{k+1}_{i} \in \Anc_G(\{v^{k+1}_{i-1}\} \cup S)$, which is what we had to show.
%          Then $\mu(B_k(i),B_{k}(i+1))$ is out of $B_{k}(i+1)$.
%          By the induction hypothesis, $B_{k}(i+1) \in \Anc_G(\{B_k(i)\} \cup S)$.
%          \Joris{NEW PROOF: since $B_k(i+1) \notin \Anc_G(B_k(i))$, we conclude $B_k(i+1) \in \Anc_G(S)$. Therefore, $B_{k+1}(i) \in \Anc_G(\{B_{k+1}(i-1)\} \cup S)$. DONE WITH THIS CASE}
%          We also assumed $B_k(i) \in \Anc_G(B_k(i+1))$.
%          Suppose $B_k(i+1) \in \Anc_G(B_k(i))$.
%          Then both nodes (and all intermediate nodes on $\mu$, by assumption) lie in the same strongly-connected component of $G$.
%          $B_k(i+1)$ is a blockable non-collider on $\mu$ by construction, and hence on $\mu$ it must point to at least one node in another strongly-connected component.
%          Therefore, $\mu(B_k(i+1),B_k(i+2))$ must be out of $B_k(i+1)$.
%          By the induction hypothesis, $B_k(i+1) \in \Anc_G(\{B_k(i+2)\} \cup S)$.
%          Hence $B_{k+1}(i) \in \Anc_G(\{B_{k+1}(i+1)\} \cup S)$.
%          This conclusion also follows (trivially) if $B_k(i+1) \notin \Anc_G(B_k(i))$ because then $B_k(i+1) \in \Anc_G(S)$.
%          \Joris{HUH? Wrong conclusion!
%          We want to conclude that $B_{k+1}(i) \in \Anc_G(\{B_{k+1}(i-1)\} \cup S)$.
%          $\mu$ must look like:
%          $$\footnotesize\begin{array}{lllllllll}
%            B_k(1)     & \dots & B_k(i-2)     & B_k(i-1)     & \suh B_k(i) & \hut \dots \hut B_k(i+1) \tuh \dots  & B_k(i+2)     & \dots & B_k(m_k) \\ 
%            B_{k+1}(1) & \dots & B_{k+1}(i-2) & B_{k+1}(i-1) & \dots  & \hut \dots \hut B_{k+1}(i) \tuh \dots & B_{k+1}(i+1) & \dots & B_{k+1}(m_{k+1})
%          \end{array}$$
%          Not so easy. 
%          Perhaps take a shortcut?
%          Or switch around the directed path between $B_k(i)$ and $B_k(i+1)$?
%          Note that $\mu(B_k(i-1),B_k(i))$ is an inducing path given $S$ that is into $B_k(i)$.
%          By switching the direction of the directed path between $B_k(i)$ and $B_k(i+1)$ we get $B_k(i-1) \dots \suh B_k(i) \tuh \dots \tuh B_k(i+1)$.
%          This is inducing (because $B_k(i)$ and all the nodes leading up to $B_k(i+1)$ are unblockable non-colliders) given $S$, and into $B_k(i+1)$.
%          Let $r_{s-1}$ denote the left neighbor of $B_k(i)$ on $\mu$.
%          Then $r_{s-1} \suh B_k(i) \hut \dots \hut B_{k}(i+1)$ is a subpath of $\mu$.
%          This means that we can take the shortcut $r_{s-1} \suh B_{k}(i+1)$.
%          Let $B_k(i-1) \sus r_1 \sus r_2 \sus \dots \sus r_{s-1} \suh B_k(i) \hut \dots \hut B_k(i+1)$ be subpath of $\mu$.
%          The $r$'s are either unblockable non-colliders, or colliders that are in $\Anc_G(S)$ or ancestor of one of their (former) neighbors.
%          If $r_{s-1}$ would be a non-collider, it would be unblockable, giving a contradiction (because then we have $r_{s-1}$ and $B_k(i+1)$ in the same scc, but $\mu(r_{s-1},B_k(i+1))$ is not directed).
%          Hence, $r_{s-1}$ must be a collider if it is not $B_k(i-1)$:
%          $B_k(i-1) \sus r_1 \sus r_2 \sus \dots \suh r_{s-1} \huh B_k(i) \hut \dots \hut B_k(i+1)$.
%          Example: $i \tuh j \tuh l \hut k$ with $j \hut l$.
%          $j$ is an unblockable non-collider so already left out.
%          }
%      \end{itemize}
  \end{itemize}
  The case that $v^k_{i} \in \Anc_G(v^k_{i-1})$ proceeds in a similar way.

  We will now show that $\pi_K$ in $H$ is $m$-open given $Z$.
  Consider a node $v^K_i$ on $\pi_K$.
%  It follows that if $\mu(v^K_{i},v^K_{i+1})$ is out of $v^K_{i}$, then $v^K_{i} \tus v^K_{i+1}$ on $\pi^K$,
%  and if $\mu(v^K_{i},v^K_{i+1})$ is out of $v^K_{i+1}$, then $v^K_{i} \sut v^K_{i+1}$ on $\pi^K$.
%  Therefore, if $v^K_{i}$ is a non-collider on $\mu$, then $v^K_{i}$ is a non-collider on $\pi^K$ as well (and is not in $Z$).
  \begin{itemize}
    \item
  Suppose $v^K_{i}$ is a non-collider on $\mu$. 
  Then it must be blockable, and hence it must be acyclically left-out-of or acyclically right-out-of.
  By the property just shown, we obtain that on $\pi^K$, we have $v^K_{i-1} \sut v^K_i$, or $v^K_{i} \tus v^K_{i+1}$, respectively. 
  Hence $v^K_{i}$ is a non-collider on $\pi^K$ as well.
  It is not in $Z$ (because all blockable non-colliders on $\mu$ are not in $Z$).
    \item
  If $v^K_{i}$ is a collider on $\mu$, then it cannot be in $\Anc_G(v^K_{i-1},v^K_{i+1})$ (otherwise the algorithm would have continued to remove $v^K_{i}$).
  Being a collider on $\mu$ that was not skipped initially, it is in $\Anc_G(Z) \setminus \Anc_G(S)$.
  %Hence it must be a collider on $\pi_K$ as well, and it must be in $\Anc_H(Z)$ \Joris{little lemma! $\Anc_G(Z) \setminus \Anc_G(S) \subseteq \Anc_H(Z)$: if $a \in \Anc_G(Z) \sm \Anc_G(S)$ then there is a directed path $v_1 \tuh \dots \tuh v_n$ with $v_1 = a$ and $v_n \in Z$ such that no node on that directed path is in $\Anc_G(S)$. Hence there is a corresponding path $v_1 \dots v_n$ in $H$ with the same nodes (but possibly different edges. For each $i<n$, $v_i \in \Anc_G(v_{i+1})$. For each $i<n$, $v_{i+1} \notin \Anc_G(v_i) \cup \Anc_G(S)$ if $G$ is acyclic. Hence the same directed path $v_1 \dots v_n$ is in $H$ if $G$ is acyclic. Aha. But since $G$ can be cyclic, we could only conclude at best that $a \in \Ant_H(Z)$...}.
  Hence it must be a collider on $\pi^K$ as well.
  Since it is in $\Anc_G(Z) \sm \Anc_G(S)$, it is in $\Ant_H(Z)$.
  \end{itemize}

  The only thing left to prove is that $\pi^K$ does not contain a subpath of the form $v^K_{i-1} \suh v^K_{i} \tut v^K_{i+1}$ or $v^K_{i-1} \tut v^K_{i} \hus v^K_{i+1}$.
  Suppose it contained a subpath of the form $v^K_{i-1} \suh v^K_{i} \tut v^K_{i+1}$.
  Then $v^K_{i}, v^K_{i+1}$ must be in the same strongly-connected component of $G$, and both are not in $\Anc_G(S)$.
  Note that $v^K_i$ can neither be acyclically left-out-of on $\mu$, nor an unblockable non-collider on $\mu$.
  %Therefore $\mu$ must look like $\suh v^K_{i} \tuh \dots \tuh v^K_{i+1}$ or $\suh v^K_{i} \hut \dots \hut v^K_{i+1}$.
  %Then $v^K_{i}$ cannot be acyclically left-out-of on $\mu$, otherwise $v^K_{i} \in \Anc_G(\{v^K_{i-1}\} \cup S)$ which would contradict the arrowhead at $v^K_{i}$.
  %\Joris{Even more, $\mu$ cannot be of the form $\dots \sut v^K_i \dots$ because then it would be an unblockable non-collider!}
  %Hence, $\mu(v^K_{i}, v^K_{i+1})$ is a directed path within a single strongly-connected component of $G$ (by assumption on the form of $\mu$).
  %In the former case, $v^K_{i}$ would be an unblockable non-collider on $\mu$, a contradiction.
  Therefore $\mu$ must look like $\suh v^K_{i} \hut \dots \hut v^K_{i+1}$.
  But then $v^K_{i}$ would have been removed by the algorithm, a contradiction.
  Hence $\pi^K$ cannot contain such a subpath $v^K_{i-1} \suh v^K_{i} \tut v^K_{i+1}$.
  One can show in a similar way that it cannot contain a subpath $v^K_{i-1} \tut v^K_{i} \hus v^K_{i+1}$ either.
  
  We now have shown the existence of a path $\pi_K$ in $H$ between $a$ and $b$ that almost qualifies for being $m$-open given $Z$, except that some of its colliders may not be in $\Anc_H(Z)$ (yet all colliders are in $\Ant^H(Z)$).
  By using Lemma~\ref{lem:sigmamag_collider_anterior_or_ancestral} repetiviely, we can construct from this an $m$-open path $\tilde{\pi}_K$ in $H$ between $a$ and $b$.
\end{proof}
We can now also define:
\begin{Def}[$im$-separation]\label{def:sigmamag_separation}
  Let $G$ be a CDMG with input nodes $J$ and output nodes $V^+ = V \dcup S$ and $H$ a $\sigma$-MAG that represents $G$ given $S$.
  Let $A,B,C \ins J \cup V$ be (not necessarily disjoint) subsets of nodes in $H$.
  We then say that:
  \begin{enumerate}
    \item \emph{$A$ is $im$-separated from $B$ given $C$ in $H$}, in symbols:
    \[ A \imPerp_H B \given C, \]
    if every walk from a node in $A$ to a node in $J \cup B$ (sic!) is $m$-blocked by $C$.
    \item If that property does not hold we will write:
      \[ A \nimPerp_H B \given C. \]
    \item We also define the special case:
      \[ A \imPerp_H B \qquad :\iff \qquad A \imPerp_H B \given \emptyset.\]
    \end{enumerate}
\end{Def}
Then we conclude:
\begin{Thm}\label{thm:sigmamag_preserves_independence_model}
  Let $G$ be a CDMG with input nodes $J$ and output nodes $V^+ = V \dcup S$ and $H$ a $\sigma$-MAG that represents $G$ given $S$.
  Let $A,B,C \ins J \cup V$ be (not necessarily disjoint) subsets of nodes in $H$.
  Then:
    \[ A \imPerp_H B \given C \iff A \isPerp_G B \given C \cup S. \]
\end{Thm}
\begin{proof}
  This follows from Lemma~\ref{lem:sigmamag_open_path_implies_cdmg_sigma-open_path} and Lemma~\ref{lem:cdmg_sigma-open_path_implies_sigmamag_open_path}.
\end{proof}

\Joris{Cannot we also define $m$-blockedness as follows?}
\begin{Def}\label{def:sigmamag_blocking2}
  Let $G$ be a CDMG with input nodes $J$ and output nodes $V^+ = V \dcup S$ and $H$ a $\sigma$-MAG that represents $G$ given $S$.
  Let $\pi$ be a walk in $H$:
  \[ \pi =\lp  v_0 \sus \cdots \sus v_n \rp.  \]
  Let $C \subseteq J \cup V$. We say that the walk $\pi$ is:
  \begin{enumerate}
    \item \emph{$C$-$m$-open} (or \emph{$m$-open given $C$}) if and only if:
      \begin{enumerate}[label=\roman*)]
        \item all colliders $v_k$ on $\pi$ are in $\Anc^H(C)$, \emph{and}
        %\item all colliders $v_k$ on $\pi$ are in $\Ant^H(C)$, \emph{and}
        \item all non-colliders $v_k$ on $\pi$ are \emph{not} in $C$, \emph{and}:
        \item it contains no subwalks of the form $v_{k-1} \suh v_k \tut v_{k+1}$ or $v_{k-1} \tut v_k \hus v_{k+1}$.
      \end{enumerate}
    \item \emph{$C$-$m$-blocked} (or \emph{$m$-blocked given $C$}) if and only if:
      \begin{enumerate}[label=\roman*)]
        \item there exists a collider $v_k$ on $\pi$ that is \emph{not} in $\Anc^H(C)$
        %\item there exists a collider $v_k$ on $\pi$ that is \emph{not} in $\Ant^H(C)$, \emph{or}:
        \item there exists a non-collider $v_k$ on $\pi$ in $C$, \emph{or}:
        \item there exists a subwalk of the form $v_{k-1} \suh v_k \tut v_{k+1}$ or $v_{k-1} \tut v_k \hus v_{k+1}$.
      \end{enumerate}  
  \end{enumerate}
  Hence the walk is $C$-$m$-open if and only if it is not $C$-$m$-blocked.
\end{Def}
So we have reverted $\Ant$ back to $\Anc$.
\begin{Prp}
  Let $G$ be a CDMG with input nodes $J$ and output nodes $V^+ = V \dcup S$ and $H$ a $\sigma$-MAG that represents $G$ given $S$.
  Let $C \subseteq J \cup V$ and $a, b \in J \cup V$.
  If there exists a walk in $H$ between $a$ and $b$ that is $m$-open given $C$, then there exists a walk in $H$ between $a$ and $b$ such that
    \begin{enumerate}[label=\roman*)]
      \item all colliders $v_k$ on $\pi$ are in $\Anc^H(C)$, \emph{and}
      \item all non-colliders $v_k$ on $\pi$ are \emph{not} in $C$, \emph{and}:
      \item it contains no subwalks of the form $v_{k-1} \suh v_k \tut v_{k+1}$ or $v_{k-1} \tut v_k \hus v_{k+1}$.
    \end{enumerate}
\end{Prp}
\begin{proof}
  Let $\pi = v_1 \sus \dots \sus v_n$ be a walk in $H$ between $v_1=a$ and $v_n=b$ that is $m$-open given $C$, that is, such that
    \begin{enumerate}[label=\roman*)]
      \item all colliders $v_k$ on $\pi$ are in $\Ant^H(C)$, \emph{and}
      \item all non-colliders $v_k$ on $\pi$ are \emph{not} in $C$, \emph{and}:
      \item it contains no subwalks of the form $v_{k-1} \suh v_k \tut v_{k+1}$ or $v_{k-1} \tut v_k \hus v_{k+1}$.
    \end{enumerate}
  Let $k_1, \dots, k_l$ denote (in ascending order) the indices of collider nodes on $\pi$. 
  We go through them one by one in ascending order and, if necessary, make a local modification of the walk $\pi$ by applying Lemma~\ref{lem:sigmamag_collider_anterior_or_ancestral} to each collider.
  We then obtain a walk that satisfies the claimed property.
\end{proof}
Note that the same applies also to an inducing walk: if there exists an `$m$-inducing' walk (with all colliders anterior to an end node) then there exists an inducing walk (with all colliders ancestor of an end node) and vice versa.
Also we need not consider $a \tuh b \tut c$ to be `$m$-inducing' in the same way that we do not need to consider it to be `$m$-open': we can instead shift attention to the path $a \tuh c$.
\begin{Def}\label{def:sigmamag-inducing_path}
  Let $G$ be a CDMG with input nodes $J$ and output nodes $V^+ = V \dcup S$ and $H$ a $\sigma$-MAG that represents $G$ given $S$.
  A walk $\pi$ in $H$ between $i$ and $j$ is called \emph{inducing} if each collider on $\pi$ is in $\Anc^H(\{i,j\})$, and it contains no non-colliders except for the endpoints. 
  If it is a path, it is called an \emph{inducing path} between $i$ and $j$.
\end{Def}
\begin{Lem}\label{lem:sigmamag-inducing_path_anterior}
  Let $G$ be a CDMG with input nodes $J$ and output nodes $V^+ = V \dcup S$ and $H$ a $\sigma$-MAG that represents $G$ given $S$.
  Let $i,j \in V$.
  There exists an inducing walk $\pi$ between $i$ and $j$ in $H$ if and only if there exists a walk $\tilde{\pi}$ in $H$ between $i$ and $j$ such that each collider on $\pi$ is in $\Ant^H(\{i,j\})$, and it contains no non-colliders except for the endpoints; $\pi$ is into $i$ iff $\tilde{\pi}$ is into $i$, and $\pi$ is into $j$ iff $\tilde{\pi}$ is into $j$.
\end{Lem}
\begin{proof}
  We go through the colliders on $\pi$ one by one and, if necessary, make a local modification of the walk by applying Lemma~\ref{lem:sigmamag_collider_anterior_or_ancestral} to each collider which is not already in $\Ant^H(\{i,j\})$.
\end{proof}


\subsubsection{What about $im$-separation in $\sigma$-PAGs?}

\Joris{Zhang (2008) writes: "It is not hard to see that if there is a definite
m-connecting path between X and Y given Z in a PAG, then in every MAG
represented by the PAG, the corresponding path is an m-connecting path between
X and Y given Z. ... A quite surprising result is that if there is an m-connecting
path between X and Y given Z in a MAG, then there must be a definite
m-connecting path (not necessarily the same path) between X and Y given Z in
its CPAG ...."
So we can just look at the presence of a definite m-connecting path in the CPAG!

The preservation of the independence model in MAGs also follows from Theorem 4.18 in Richardson \& Spirtes (2002).
Lemmas 20 and 26 in Zhang (2008) are useful! The latter points to Lemma 5.1.7 in Zhang (2006), the proof of which is on p.~208--216, the former points to Lemma 5.1.3.  Note that it would be VERY useful if we could extend (the first part of) Lemma 5.1.7 to the selection bias setting!
This also provides us with what we need to extend the do-calculus for MAGs to a do-calculus for PAGs.
Note that the negation of `definitely m-connecting' is `possibly m-closed'.}
